%% 使用 scutthesis 文档类生成南京大学学位论文的示例文档
%%
%% 作者：胡海星，starfish (at) gmail (dot) com
%% 项目主页: http://haixing-hu.github.io/nju-thesis/
%%
%% 本样例文档中用到了吕琦同学的博士论文的提高和部分内容，在此对他表示感谢。
%%
\documentclass[phd,nobackinfo]{scutthesis}
%% scutthesis 文档类的可选参数有：
%%   nobackinfo 取消封二页导师签名信息。注意，按照南大的规定，是需要签名页的。华工则不需要，
%%   故开启该选项
%%   phd/master/bachelor 选择博士/硕士/学士论文

% 使用 blindtext 宏包自动生成章节文字
% 这仅仅是用于生成样例文档，正式论文中一般用不到该宏包
\usepackage[math]{blindtext}

%%%%%%%%%%%%%%%%%%%%%%%%%%%%%%%%%%%%%%%%%%%%%%%%%%%%%%%%%%%%%%%%%%%%%%%%%%%%%%%
% 设置《国家图书馆封面》的内容，仅博士论文才需要填写

% 设置论文按照《中国图书资料分类法》的分类编号
\classification{0175.2}
% 学校代号，华工代号为10561
\schoolid{10561}
% 学号
\studentnum{201710105133}
% 论文的密级。需按照GB/T 7156-2003标准进行设置。预定义的值包括：
% - \openlevel，表示公开级：此级别的文献可在国内外发行和交换。
% - \controllevel，表示限制级：此级别的文献内容不涉及国家秘密，但在一定时间内
%   限制其交流和使用范围。
% - \confidentiallevel，表示秘密级：此级别的文献内容涉及一般国家秘密。
% - \clasifiedlevel，表示机密级：此级别的文献内容涉及重要的国家秘密 。
% - \mostconfidentiallevel，表示绝密级：此级别的文献内容涉及最重要的国家秘密。
% 此属性可选
% 华工格式与标准略为不同，公开级为空。
\securitylevel{\openlevel} % 非保密论文此处值为空，否则格式为密级加年限
% 如下
% \securitylevel{\confidentiallevel \bigstar 5年}
% 设置论文按照《国际十进分类法UDC》的分类编号
% 该编号可在下述网址查询：http://www.udcc.org/udcsummary/php/index.php?lang=chi
% \udc{004.72}
% 国家图书馆封面上的论文标题第一行，不可换行。此属性可选，默认值为通过\title设置的标题。
\nlctitlea{二维硼化钛的高通量计算}
% 国家图书馆封面上的论文标题第二行，不可换行。此属性可选，默认值为空白。
\nlctitleb{及其超导性质的第一性原理研究}
% 国家图书馆封面上的论文标题第三行，不可换行。此属性可选，默认值为空白。
% \nlctitlec{（第三行提名或副提名）}
% 导师的单位名称及地址
\supervisorinfo{华南理工大学物理与光电学院~~广州市天河区五山路381号~~510641}
% 答辩委员会主席
\chairman{张三丰~~教授}
% 第一位评阅人
\reviewera{阳顶天~~教授}
% 第二位评阅人
\reviewerb{张无忌~~副教授}
% 第三位评阅人
\reviewerc{黄裳~~教授}
% 第四位评阅人
\reviewerd{郭靖~~研究员}

%%%%%%%%%%%%%%%%%%%%%%%%%%%%%%%%%%%%%%%%%%%%%%%%%%%%%%%%%%%%%%%%%%%%%%%%%%%%%%%
% 设置论文的中文封面
\titlea{二维硼化钛的高通量计算}
\titleb{及其超导性质的第一性原理研究}

% 论文作者姓名
\author{余桔颂}
% 论文作者联系电话
\telphone{18933996792}
% 论文作者电子邮件地址
\email{morty.yeu@gmail.com}
% 论文作者学生证号
% \studentnum{MG0033011}
% 论文作者入学年份（年级）
\grade{2017}
% 导师姓名职称
\supervisor{杨小宝~~教授}
% 导师的联系电话
\supervisortelphone{???}
% 论文作者的学科与专业方向
\major{物理电子学}
% 论文作者的研究方向
\researchfield{计算凝聚态物理}
% 论文作者所在院系的中文名称
\department{物理与光电学院}
% 论文作者所在学校或机构的名称。此属性可选，默认值为``南京大学''。
\institute{华南理工大学}
% 论文的提交日期，需设置年、月、日。
\submitdate{2021年4月10日}
% 论文的答辩日期，需设置年、月、日。
\defenddate{2021年6月1日}
% 学位授予日期，需设置年、月、日。
\awarddate{2021年6月20日}
% 论文的定稿日期，需设置年、月、日。此属性可选，默认值为最后一次编译时的日期，精确到日。
%% \date{2013年5月1日}

%%%%%%%%%%%%%%%%%%%%%%%%%%%%%%%%%%%%%%%%%%%%%%%%%%%%%%%%%%%%%%%%%%%%%%%%%%%%%%%
% 设置论文的英文封面

% 论文的英文标题，不可换行
\englishtitle{A Research on Network Infrastructures for Data Centers}
% 论文作者姓名的拼音
\englishauthor{WEI Xiao-Bao}
% 导师姓名职称的英文
\englishsupervisor{Professor CHEN Jin-Nan}
% 论文作者学科与专业的英文名
% \englishmajor{Computer Software and Theory}
% 论文作者所在院系的英文名称
\englishdepartment{Department of Computer Science and Technology}
% 论文作者所在学校或机构的英文名称。此属性可选，默认值为``Nanjing University''。
\englishinstitute{South China University of Technology}
% 论文作者所在学校或机构的英文所属地，默认为``Guangzhou, China''。
\englishlocation{Guangzhou, China}
% 论文完成日期的英文形式，它将出现在英文封面下方。需设置年、月、日。日期格式使用美国的日期
% 格式，即``Month day, year''，其中``Month''为月份的英文名全称，首字母大写；``day''为
% 该月中日期的阿拉伯数字表示；``year''为年份的四位阿拉伯数字表示。此属性可选，默认值为最后
% 一次编译时的日期。
% \englishdate{May 1, 2013}

%%%%%%%%%%%%%%%%%%%%%%%%%%%%%%%%%%%%%%%%%%%%%%%%%%%%%%%%%%%%%%%%%%%%%%%%%%%%%%%
% 设置论文的中文摘要

% 设置中文摘要页面的论文标题及副标题的第一行。
% 此属性可选，其默认值为使用|\title|命令所设置的论文标题
% \abstracttitlea{数据中心网络模型研究}
% 设置中文摘要页面的论文标题及副标题的第二行。
% 此属性可选，其默认值为空白
% \abstracttitleb{}

%%%%%%%%%%%%%%%%%%%%%%%%%%%%%%%%%%%%%%%%%%%%%%%%%%%%%%%%%%%%%%%%%%%%%%%%%%%%%%%
% 设置论文的英文摘要

% 设置英文摘要页面的论文标题及副标题的第一行。
% 此属性可选，其默认值为使用|\englishtitle|命令所设置的论文标题
\englishabstracttitlea{A Research on Network Infrastructures}
% 设置英文摘要页面的论文标题及副标题的第二行。
% 此属性可选，其默认值为空白
\englishabstracttitleb{for Data Centers}

%%%%%%%%%%%%%%%%%%%%%%%%%%%%%%%%%%%%%%%%%%%%%%%%%%%%%%%%%%%%%%%%%%%%%%%%%%%%%%%
\begin{document}

%%%%%%%%%%%%%%%%%%%%%%%%%%%%%%%%%%%%%%%%%%%%%%%%%%%%%%%%%%%%%%%%%%%%%%%%%%%%%%%

% 制作国家图书馆封面（博士学位论文才需要）
\makenlctitle
% 制作中文封面
\maketitle
% 制作英文封面
\makeenglishtitle

\makelicense


%%%%%%%%%%%%%%%%%%%%%%%%%%%%%%%%%%%%%%%%%%%%%%%%%%%%%%%%%%%%%%%%%%%%%%%%%%%%%%%
% 开始前言部分
\frontmatter

%%%%%%%%%%%%%%%%%%%%%%%%%%%%%%%%%%%%%%%%%%%%%%%%%%%%%%%%%%%%%%%%%%%%%%%%%%%%%%%
% 论文的中文摘要
\begin{abstract}
复杂网络的研究可上溯到20世纪60年代对ER网络的研究。90年后代随着Internet
的发展，以及对人类社会、通信网络、生物网络、社交网络等各领域研究的深入，
发现了小世界网络和无尺度现象等普适现象与方法。对复杂网络的定性定量的科
学理解和分析，已成为如今网络时代科学研究的一个重点课题。

在此背景下，由于云计算时代的到来，本文针对面向云计算的数据中心网络基础
设施设计中的若干问题，进行了几方面的研究。………………
% 中文关键词。关键词之间用中文全角分号隔开，末尾无标点符号。
\keywords{小世界理论；网络模型；数据中心}
\end{abstract}

%%%%%%%%%%%%%%%%%%%%%%%%%%%%%%%%%%%%%%%%%%%%%%%%%%%%%%%%%%%%%%%%%%%%%%%%%%%%%%%
% 论文的英文摘要
\begin{englishabstract}
\blindtext
% 英文关键词。关键词之间用英文半角逗号隔开，末尾无符号。
\englishkeywords{Small World, Network Model, Data Center}
\end{englishabstract}

% %%%%%%%%%%%%%%%%%%%%%%%%%%%%%%%%%%%%%%%%%%%%%%%%%%%%%%%%%%%%%%%%%%%%%%%%%%%%%%%
% % 论文的前言，应放在目录之前，中英文摘要之后
% %
% \begin{preface}
%
% 复杂网络的研究可上溯到20世纪60年代对ER网络的研究。90年后代随着Internet
% 的发展，以及对人类社会、通信网络、生物网络、社交网络等各领域研究的深入，
% 发现了小世界网络和无尺度现象等普适现象与方法。对复杂网络的定性定量的科
% 学理解和分析，已成为如今网络时代科学研究的一个重点课题。
%
% 在此背景下，由于云计算时代的到来，本文针对面向云计算的数据中心网络基础
% 设施设计中的若干问题，进行了几方面的研究。本文的创造性研究成果主要如下
% 几方面：
%
% \begin{enumerate}
% \item 基于簇划分的思想，提出并设计了WarpNet网络模型。该网络模型基于随机
%   散列，以节点微路由链接多种散列分布，实现网络互联。并对网络的带宽等指
%   标进行理论分析并给出定量描述。最后对比了理论分析、仿真测试结果，并在
%   实际物理环境中进系真实部署，通过6节点的小规模实验以及1000节点虚拟机的
%   大规模实验，表明该模型的理论分析、仿真测试与实际实验吻合，并在网络性
%   能、容错能力、伸缩性灵活性方面得到较大提升。
% \item 提出DS小世界模型并构造SIDN网络，解决了把小世界理论应用于数据中心
%   网络布局构建中的最大度限制问题。分析了在带有最大度限制约束下，所构成
%   网络的平均直径、网络总带宽、对故障的容错能力等各项网络参数。理论分析
%   与仿真实验证明，SIDN网络具有很好的扩展能力，网络总带宽与网络规模成近
%   似线性增长的关系；具有很强的容错能力，链路损坏与节点损坏几乎无法破坏
%   网络的联通性，故障率对网络性能的影响与破坏节点/链路占总资源比率线性相
%   关。
% \item 分析了无尺度网络在数据中心网络构建应用中的理论方面问题。在引入节
%   点最大度限制之后，给出无尺度网络的各项网络参数。并进一步分析了交换机
%   节点以及计算节点两种角色在不同比率的组合下对网络性能的影响，给出最高
%   性价比的比率参数。最后通过理论分析与仿真实验证明，在引入了无尺度现象
%   之后，提高了网络的聚类系数，从而显著的提升了网络的性能。
%
% \item 针对网络模型研究这一类工作的共性，设计构造通用验证平台系统。以海
%   量虚拟机和虚拟分布式交换机的形式，实现了基于少量物理节点，对大规模节
%   点的模拟。其模拟运行的过程与真实运行在实现层面完全一致，运行的结果与
%   真实环境线性相关。除为本文所涉若干网络模型提供验证外，可进一步推广到
%   更为广泛的领域，为各种网络模型及路由算法的研究工作，提供分析、指导与
%   验证。
% \end{enumerate}
%
% 复杂网络的研究可上溯到20世纪60年代对ER网络的研究。90年后代随着Internet
% 的发展，以及对人类社会、通信网络、生物网络、社交网络等各领域研究的深入，
% 发现了小世界网络和无尺度现象等普适现象与方法。对复杂网络的定性定量的科
% 学理解和分析，已成为如今网络时代科学研究的一个重点课题。
%
% 在此背景下，由于云计算时代的到来，本文针对面向云计算的数据中心网络基础
% 设施设计中的若干问题，进行了几方面的研究。本文的创造性研究成果主要如下
% 几方面：
%
% \begin{enumerate}
% \item 基于簇划分的思想，提出并设计了WarpNet网络模型。该网络模型基于随机
%   散列，以节点微路由链接多种散列分布，实现网络互联。并对网络的带宽等指
%   标进行理论分析并给出定量描述。最后对比了理论分析、仿真测试结果，并在
%   实际物理环境中进系真实部署，通过6节点的小规模实验以及1000节点虚拟机的
%   大规模实验，表明该模型的理论分析、仿真测试与实际实验吻合，并在网络性
%   能、容错能力、伸缩性灵活性方面得到较大提升。
% \item 提出DS小世界模型并构造SIDN网络，解决了把小世界理论应用于数据中心
%   网络布局构建中的最大度限制问题。分析了在带有最大度限制约束下，所构成
%   网络的平均直径、网络总带宽、对故障的容错能力等各项网络参数。理论分析
%   与仿真实验证明，SIDN网络具有很好的扩展能力，网络总带宽与网络规模成近
%   似线性增长的关系；具有很强的容错能力，链路损坏与节点损坏几乎无法破坏
%   网络的联通性，故障率对网络性能的影响与破坏节点/链路占总资源比率线性相
%   关。
% \item 分析了无尺度网络在数据中心网络构建应用中的理论方面问题。在引入节
%   点最大度限制之后，给出无尺度网络的各项网络参数。并进一步分析了交换机
%   节点以及计算节点两种角色在不同比率的组合下对网络性能的影响，给出最高
%   性价比的比率参数。最后通过理论分析与仿真实验证明，在引入了无尺度现象
%   之后，提高了网络的聚类系数，从而显著的提升了网络的性能。
%
% \item 针对网络模型研究这一类工作的共性，设计构造通用验证平台系统。以海
%   量虚拟机和虚拟分布式交换机的形式，实现了基于少量物理节点，对大规模节
%   点的模拟。其模拟运行的过程与真实运行在实现层面完全一致，运行的结果与
%   真实环境线性相关。除为本文所涉若干网络模型提供验证外，可进一步推广到
%   更为广泛的领域，为各种网络模型及路由算法的研究工作，提供分析、指导与
%   验证。
% \end{enumerate}
%
% \vspace{1cm}
% \begin{flushright}
% 韦小宝\\
% 2013年夏于南京大学
% \end{flushright}
%
% \end{preface}

%%%%%%%%%%%%%%%%%%%%%%%%%%%%%%%%%%%%%%%%%%%%%%%%%%%%%%%%%%%%%%%%%%%%%%%%%%%%%%%
% 生成论文目次
\tableofcontents

%%%%%%%%%%%%%%%%%%%%%%%%%%%%%%%%%%%%%%%%%%%%%%%%%%%%%%%%%%%%%%%%%%%%%%%%%%%%%%%
% 生成插图清单。如无需插图清单则可注释掉下述语句。
\listoffigures

%%%%%%%%%%%%%%%%%%%%%%%%%%%%%%%%%%%%%%%%%%%%%%%%%%%%%%%%%%%%%%%%%%%%%%%%%%%%%%%
% 生成附表清单。如无需附表清单则可注释掉下述语句。
\listoftables

%%%%%%%%%%%%%%%%%%%%%%%%%%%%%%%%%%%%%%%%%%%%%%%%%%%%%%%%%%%%%%%%%%%%%%%%%%%%%%%
% 开始正文部分
\mainmatter

%%%%%%%%%%%%%%%%%%%%%%%%%%%%%%%%%%%%%%%%%%%%%%%%%%%%%%%%%%%%%%%%%%%%%%%%%%%%%%%
% 学位论文的正文应以《绪论》作为第一章
\chapter{绪论}\label{chapter_introduction}
\section{二维过渡金属硼化物研究进展}

硼元素作为周期表中碳元素的近邻元素有三个价层电子，由于其不能够通过经典的二中心二电子硼-硼键形成全满的八电子构型，硼元素在形成化合物时会展现出极其丰富的成键形式，这个原因导致了硼的化合物有极为丰富的结构多样性。稳定的硼材料中通常都存在缺点子的成键类型，比如，在许多硼材料中都能够发现三中心二电子键的存在。其中，二维硼平面的多样性丰富且存在许多新奇的物理化学性质。近些年，基于硼元素的二维材料设计吸引了广泛的研究关注。大量的二维稳定硼结构被理论预测报道，其中一些还展现出非常优秀的电学、力学性质，如能带中存在狄拉克锥、力学上有高的延展性等。根据理论预测硼元素在部分金属衬底上更易成核并生长为二维材料[]，三个不同的研究组分别成功在不同衬底上合成了二维硼材料[][][]。有趣的是，二维硼在银衬底上表现出金属特性[]，而在铜衬底则表现半导体性质[]。

过渡金属因为外层具有未满的价层d轨道，其形成化合物时存在许多可变的价态，这个原因同样也导致了含有过渡金属的化合物会表现许多丰富的物理化学性质。硼和金属形成化合物，如在硼平面或硼团簇中嵌入金属元素，尤其是过渡金属元素，使得硼元素化合物在结构多样性方面更加的丰富，同时也可能产生更为丰富的物理化学性质。在这一节我们将介绍二维硼材料的多样性以及在二维硼中嵌入过渡金属元素后所得到的新颖化合物及其性质的理论计算方面的研究。

\subsection{二维硼材料}
石墨烯[]的发现，使得二维材料如硼氮平面[]，硅烯[]，锗烯[]，磷烯[]，过渡金属双硫化合物[]，砷烯[]等，在近十多年间得到了广泛的研究关注，越来越多的二维材料被实验合成或被理论发现。由于二维材料独特的物理化学性质（如费米能级附近的高导电性，高导热性，高机械强度），二维材料在电器件、能源储存等领域有重要的应用前景。

从2015年第一次在银衬底上合成硼平面至今，二维硼的研究在凝聚态物理、化学和纳米材料领域都有长足的发展。由于其独特的物理化学性质，二维硼材料在应用上有很大的潜质。

硼烯是由硼原子所构成的单层结构，它能够在高真空环境下生长在银衬底上[]。所有被实验合成四个相（2Pmm，beta12,chi3,蜂窝相）均为金属相。性质方面，有研究分别报道了其机械性能，电子结构，晶格热导，超导，原子吸附等[]，我们将在？小节中详细介绍。由于硼平面的能量相比于体相的硼高出了不少，因此硼平面的合成难度很大。其中2-Pmmm相的硼平面的能量和alpha-六方晶系的体相硼之间的能量差为555meV/atom []。四个已经被合成的二维硼的相均是在实验合成之前就已经被实验所报道，这展现了理论预测的巨大作用。扫描隧道显微镜（STM）获取的的图像和结构的建模如图？所示，其中硼烯2-Pmmn相显示很大的各项电子异性，同时实验还报道了它特殊的机械性质。该相在衬底上呈褶皱分布，隆起处垂直方向的高度为0.91A。beta12和chi3相可以在银的111面衬底上合成，这两相都是平面结构，竖直方向上没有隆起，这两个相中水平方向上的硼原子的分布是不同的。在2019年于铝的111面衬底上生长得到的石墨烯形状的硼烯结构，该结构在铝111面衬底上比在银111面衬底上要稳定的多。原因是，与石墨烯相比，该硼烯中每个硼原子有一个电子的缺失，但银的111面相比铝的衬底不容易出现电荷转移，而铝恰好能向硼平面转移电子以稳定结构。这四个稳定的二维硼烯的合成大大激励了研究人员对硼烯的研究热度，同时也为硼烯在纳米器件和能源领域的应用提供了可能。

其余的许多二维硼的同素异形体都是由理论预测所报道的[]。硼烯的稳定性与硼平面中空位的数量有重要的联系，我们设平面三角网格中六角空位的比例为x，则x可能的取值必定在完全的三角网格和石墨烯形的六角格子的空位率之间。如图？[]所示，硼平面的形成能的空位率的变化呈现V型的关系，即从完全的三角格子向六角格子过渡，形成能先随着空位率的增加而降低（变得更稳定），在达到一定数值后又随着空位率的增加而升高（变得不稳定）。在空位率为x=1/9时，可以得到对应的最为稳定的结构，同时在相近的空位率为x=1/8和x=2/15时，同样有很低的形成能，此类结构都有可能稳定地存在。与其他二维材料如石墨烯、硅烯、锗烯、氮化硼以及黑磷所不同的是，硼烯的各种平面结构中在最稳定的结构附近，存在大量形成能相近的构型。

除了受到空位率的影响，还如前面所列举的，硼烯的稳定性还受到了衬底的影响。Yakobson等提出[]，无衬底时的基态结构虽然形成能较低，但生长在衬底时却不一定是最稳定的。如图？所示，在银的111面衬底上，beta12相是最稳定的硼烯，其空位率为1/6.第二稳定和第三稳定的构型的空位率分别为x=1/8和x=1/12.相应的形成能比最稳定的构型要高出0.4meV和2.1meV。在铜的111表面上能量最低的三种构型的空位率均为x=1/6,相比于能量最低的构型，次低和第三稳定的构型的形成能分别高出15.1meV的28.9meV[]。这个结果表明，在铜的衬底上更有可能形成特定的统一的硼烯构型。

除了空位率和衬底外，官能团修饰是调节硼烯物理和化学性质的另外一个有效手段。比如在其他二维材料如石墨烯、硅烯、锗烯中，能够使用氢化的方法来增加或打开材料的带隙，将材料从金属变为半导体[]。同时官能团修饰后的材料在结构上也会发生巨大的改变。石墨烯是完全平面的二维材料，但完全氢化后石墨烯结构上呈现隆起状。另外通过氢化还能改变平面结构的稳定性，比如对2-Pmmn相的硼烯进行氢化，声子谱在Gamma点处的虚频消失，结构上更加稳定[]。同时氢化还会改变构型的电子结构[]，在对2-Pmmn相的硼烯氢化后，原结构对应电子结构中电子占据的反键部分的电子向氢原子发生转移，使得面内的成键轨道完全占据，反键轨道无占据，从而使结构更加稳定。而且，该氢化后的结构在布里渊区的X-Gamma方向存在狄拉克锥的电子结构，狄拉克点处的电子速度为3.5x106m/s[]。另外，氢化后结构的力学性质也得到了改善，其面内两个方向的杨氏模量分别为172.24N/m和110.59N/m[]。

热传导性能是纳米器件使用性能和耐用性的重要指标。有一系列的文献[][]报道了有关这一性质的研究。硼烯2-Pmmn相为各向异性，其热导也表现出极大的各向异性。室温时，2-Pmmn相延锯齿状（zigzag）方向和延椅式状（armchair）方向的热导分别为75.9W/mK和147W/mk。这两个方向的声子平均有效自由路程为16.7nm和21.4nm。alpha硼平面具有各向同性，其热导为14.34W/mK。硼烯中热导主要是由高频声子的振动模式所贡献的，这不同于在其他二维材料如石墨烯[]，硅烯[]，磷烯[]，主要是低频的声子模式决定了热传导性质。

从实验上合成得到的单层硼平面结构，研究人员对其进行了完整的电子结构性质的研究。理论计算的结果显示，实验上得到的单层硼烯平面，原子起伏的2-Pmmn相，β12相以及χ3相，都是金属[]，Yakobson 等[]通过理论计算预测，这种金属单层硼平面应该具有超导电性质[]。同时，他们还进一步研究了该类金属二维材料其表面的光学性质[]。Feng等[]人通过角分辨光电子能谱(ARPES)测量技术揭示（如图？），在实验上对应的两种单层硼的相的电子结构能带中存在非平庸的狄拉克点，因此认为其有可能是二维拓扑材料。之后，又有进一步的理论计算工作预言了在其他类型的单层硼平面的能带结构中，不但存在各项异性的狄拉克锥， 而且还存在有节线(node-line)特征的非平庸的色散关系[]。Du等[]人的理论工作指出，六角晶格的单层硼平面，如果硼原子表面用氢原子修饰，也会实现各向异性的狄拉克点。

需要提及的是，理论预测不但辅助实验合成得到了二维单原子层硼平面，理论还预言了硼存在的二维薄膜。然而，到目前为止并没有发现硼的层状体相材料的存在，目前已知所有的硼的体相构型都具有高配位的极其复杂的成键方式。如果使用全局构型搜索的方法，当不限制搜索时二维的材料的厚度时，有理论工作表明，多层的硼结构会比孤立单层硼平面更加稳定。比如Zhou等[]人通过遗传算法理论预测了稳定的非单层的有狄拉克点的二维硼材料。Du等[]的也发现了一种新型的具有双狄拉克色散的二维硼薄膜。实验上，2015 年南京航空航天大学科研团队在铜表面以硼和B2O3为前驱体，通过加H2还原的方式得到了二维硼薄膜[]，该硼薄膜被认为是从体相硼 γ28 中切割出的单层结构。理论计算表明，该新的硼薄膜材料是一种宽直接带隙的半导体材料，其相应的带隙值为 2.07 eV，使其有可能被作为具有优秀光吸收性质的材料。然而，后续又有其他的理论工作[]认为这种二维薄膜其实是金属，但如果在其表面通过氢原子进行修饰，可以容易的实现金属和半导体之间的调控。

单层硼平面的结构多样性非常丰富，然而，大多数理论计算和实验发现的单层硼平面都是金属[]，Yakobson 等[][]在多篇文献中的认为，单层硼平面之所以呈现金属性，是因为其体系的费米能级一定会穿过由体系构型中平面结构面外垂直方向的pz 轨道所构成的能带。但是体系本身的面内轨道(s, px 和 py轨道) 在费米能级处是可能存在带隙的。所以普遍的观点认为，单层硼平面所有构型都不会表现出半导体性质。但我们课题组2017年的工作[]发现了表现为半导体的几种硼平面构型，这大大的拓展了硼平面的性质。

除此之外，硼平面的化学活性也有不少的相关研究。通过将孤立的锂原子吸附在石墨烯表面，有理论计算指出[]，通过控制锂元素的覆盖度变化，可以得到不同的超导转变温度的二维超导体。随后，有另外的研究者[]将同样的思路拓展到单层硼平面的锂吸附性质上，通过结合集团展开来评估体系能量，他们预测了一系列随覆盖度变化而变化的 Li/B 超导体，其中所预测的 Li2B7构型的超导转变温度为6.2 K。 该研究说明了，通过控制覆盖度以及平面上的孤立原子的吸附位点，可制备出一系列超过液氦温度的传统BCS二维超导体。

同时，因为硼单层平面与金属锂具有强的吸附作用，被锂修饰的硼单层平面材料可以被设计用来作为高性能的储氢材料。有理论[]预测稳定了具有 1/8 空位浓度的硼单层平面就可能具有较理想的储氢性能。并且，实验上合成得到的硼平面结构也被用来实现锂离子电池的设计。通过理论计算[]发现，2-Pmmn相的起伏三角格子的二维硼结构沿其起伏形成的通道方向，其锂离子的扩散非常容易，对应的扩散势垒仅有 2.6 meV。并且，实验上所合成得到的两种基态构型 β12相和 χ3相本身就有较多的空位的分布，空位连接成链时其可以被看作是理想的离子扩散通道[]，但不理想的是，由于硼原子的缺电子性，这些空位对应的是较强的金属吸附，导致了锂原子沿这类链状空位路径扩散时的势垒约为0.6 eV。如下图？所示，分别是 β12 相和 χ3相锂离子扩散和 钠离子扩散的模拟路径。相比于其他的二维材料，单层硼烯具有较大的吸附能。理论计算表明[]，在实验上单层 β12 平面的六角空位处吸附孤立的金属原子铋，可以作为一个非常有效的单原子催化载体从而可以实现水分解，β12平面还可作为绿色环境材料用于实现对温室气体CO2的高效捕获[]

\subsection{过渡金属硼平面及其理论计算研究进展}
研究原子团簇，及原子团簇结构和性质之间的关系，对理解化学键本质以及设计性质可调控的纳米材料有重要的意义。硼由于其缺电性，较少的硼原子在形成团簇时倾向于形成平面团簇结构[]，这一特征和硼的体相结构及体相含硼化合物中的硼复杂网络结构形成鲜明对比。硼的阴离子团簇最大可达到23个原子[]，在硼原子数量为16的硼阳离子团簇和硼原子数量为20的中性硼原子团簇中，均发现了结构从二维到三维的转变[][]。所有实验发现的硼平面团簇，都是由一圈通过二中心二电子（2c-2e）硼硼键所形成的硼环围绕一个中心原子，硼和中心原子通过sigma键连接，整个平面团簇在水平方面形成一个大pi键来稳定团簇。我们可以注意到，在阴离子硼团簇中，中心原子通常和多个硼原子形成多中心键，一个典型的例子是B19-（B@B5@B13）[]，其中就包含两种不同的离域pi键。这种离域表示了中心硼原子与中间层的B5环，B5环与最外层B13环之间的相互作用。有趣的是，内层的B@B5结构可以在B13环中自由旋转，这使其成为了分子器件的一个潜在候选。

当八个或九个硼原子成环时，可以在环的中央放入金属离子而形成稳定的团簇结构，这两类团簇是有高对称性的轮状分子，其中硼元素组成的环符合（4N+2）休克尔规则，使得结构能够稳定存在[]。实验发现报道的有AlB7-（C7v）和AlB-（C8v）[]，结构中的铝元素分别和B7与B8环通过离子相互作用结合，金属-硼之间的成键并不离域。2013年Romanescu等通过激光蒸发喷射超快团簇束的方法得到了一系列由过渡金属为中心与硼环形成的轮状分子结构，Co@B8-,Ru@B9-，Rh@B9-,Ir@B9-，Nb@B9-以及Ta@B9-[][][]，这些实验发现的结构也同样被理论证实是能量最低的平面团簇结构。

Qu等[]使用粒子群算法对钛硼比例为1：4的结构进行了搜索，找到了全平的二维单层钛硼结构，如图?所示，其中钛原子被八个硼原子环绕形成八配位的钛硼轮状结构，八个硼原子到钛原子的距离均相等。他们计算了该结构的动力学性质，认为其十分稳定因而可能够被实验所合成。同时他们认为该结构的稳定性源于钛和硼形成的八配位的稳定轮状团簇，因为钛向硼提供了电子弥补了硼的缺电性使硼原子之间形成了sp2杂化轨道构成的稳定网络。受该结构启发，研究者还对该钛硼结构进行了其他金属原子的替换，将钛替换为钒，镉，锰，钨元素形成平面或准平面结构并讨论了这些结构的稳定性。

Bo等[]通过第一性原理结合粒子群算法，发现并研究了两种Mo2B2单层的稳定性，电子结构，晶格动力学性质以及其在电池材料中作为阴极材料时的性质。声子谱和电子结构的信息展示了正方和三方晶系Mo2B2结构(如图？)的高电导性质以及高的稳定性。这两个结构对锂原子和钠原子的吸附能在电池的工作范围内都非常大，确保了其工作效率。其中正方晶系Mo2B2的锂离子和钠离子容量均为约251mAhg-1, 三方晶系Mo2B2的锂离子和钠离子容量分别为约251mAhg-1和188mAhg-1。三方晶系的Mo2B2对锂离子和钠离子的扩散势垒分别为0.023和0.013eV，四方Mo2B2对锂钠离子的扩散势垒分别为0.029和0.010eV，扩散势垒均不高说明他们有很好的充放电效率。相同研究组的研究人员还在另一工作中对六方晶系Ti2B2进行了锂离子和钠离子电池性质的研究，该材料的锂钠离子的容量分别为456mAhg-1和1027mAhg-1。值得一提的是该材料对锂离子和钠离子的扩散势垒非常低分别仅为0.017和0.008eV，因此推断其具有非常优良的充放电性质。

Tang等[]使用粒子群算法预测了钴原子和硼原子形成的单层平面化合物CoB6（结构如图？所示），并研究了该材料的磁电性质。动力学计算显示该材料有很好的稳定性，同时他们认为该平面CoB6是稳定的二维铁磁材料，而且其磁矩和磁稳定性能够通过对材料施加应力来提高。他们还研究了对该材料可能的合成方法，认为可以通过在已经被实验所报道的delta4相硼单层平面上吸附钴原子或直接通过化学生长得到。

\section{硼化物的超导性质研究进展}
原子质量小的元素所构成的金属材料会展现新颖的电子结构，硼作为原子质量较轻的一个元素，单质硼或硼化合物会出现相对较高的超导温度。2001年Eremets等[]通过金刚石对顶针在超高压160GP下将硼从非金属相转变为金属相，并表现出超导性质，当压力从175GPa升高到250GPa时，其超导温度从6K提升到了11.2K，即每GPa压力增加可提高临界温度0.05K，理论计算表明该高压材料的超导为BCS机制的超导，对临界温度的计算结果与实验相符合。同年，Nagamatsu等[]报道了具有最高临界温度的非铜氧型的超导体MgB2，其常压下超导温度为39K。与MgB2由相同结构的化合物AlB2同样表现出超导性能，但常压下的超导温度为较小，却能够通过多种手段调节使其超导温度提升。同类型结构，不同金属元素的XB2型化物也同样报道存在超导性，如ZrB2[]，NbB2+x[]以及TaB2+x[]。这种材料在结构上表现为硼原子以六角格子网络连接，金属原子以插层形式位于硼网络层之间，如图？。这种层状嵌套的结构特征启发人们，可能会存在二维的硼或硼化物和超导材料。

二维材料的超导性质在纳米超导量子器件[]和纳米超导晶体管[]中都有潜在重要应用。在石墨烯上通过沉积金属原子的方式，能够使其出现超导性。2007年首次报道了通过等离子方式在石墨烯上覆盖金属而诱发超导[]，后续的理论计算表明其超导温度为LiC6单层8.1K，和CaC6单层1.4K，通过增加石墨烯层数，超导临界温度会发生改变。2017年，Gao等[]对实验上能够在银111面上生长的beta12和chi3两种硼平面进行了超导性质的计算，预测其临界超导温度分别为18.7K和24.7K。根据2012年Profeta等[]报道的在石墨烯上附着锂元素能够得到有超导性质的层状结构，2016年Wu等[]对锂在硼平面上沉积的结构进行了超导性质的计算，发现了一系列可能稳定存在的锂硼层状超导结构。2019年Yan等[]预测了两种稳定的Mo2B2结构并计算了它们的超导临界温度为3.9K和0.2K，并发现其电声耦合主要是由于钼原子的低频振动所导致的。

\section{结构搜索与高通量方法在硼材料研究中的作用}
原子在空间中能够以各种各样的形式堆积，从而形成各种类型的凝聚态物质，物质结构，是深入理解物质物理化学性质的非常重要的信息之一。《自然》杂志主编John Maddox曾在《自然》上发表评论称，“依据且仅依据物质的化学组分来确定物质的结构，是物理学的一个重要挑战”。通过计算材料学，从理论上预测和设计具有新颖或优良性质的材料具有极其重要的意义。原因是，理论计算的成本相比于实验较低，且不乏重要的前瞻性，对实验研究中以试错为主要探索方式的方法过程具有很强的指导意义。另外，理论计算的探索和研究，能够帮助研究人员从本质来理解材料结构的其物质性质之间的内在联系，并提供通过结构改变物质性质的关键信息，指导新材料的设计和发现。这样的研究方法对各种物质性质的研究都有着重要的价值。

在过去，科学家只能对已有的材料进行研究，通过实验表征其各种物质性质。因此，对物质性质的改性也往往只能限制在已有的物质结构的基础上，这大大限制了对材料性质的探索。近年来，科学家发展了一些新的结构搜索程序比如基于粒子群优化机制的CALYPSO程序[]、基于遗传算法的USPEX程序[]。仅需要材料元素种类和配比, 不需要任何结构信息, 产生大量可能的结构, 然后对这些结构进行优化和能量比较, 得出能量比较低的结构, 再进一步研究这些结构的性质和电子结构性质。结构搜索程序的日趋成熟加上服务器计算能力的提高, 极大地推进了探索新型材料的进程. 对实验上已经存在的材料, 研究者采用结构搜索程序对其可能的结构也进行了探索和尝试 。

硼原子能够构成许多原子数不同尺寸各异的硼团簇，且其具有新奇的电子结构特征。通过CALYPSO，Wang等[]2014年报道了38个硼原子组成的硼团簇，该团簇有很好的稳定性，较大的能带和高的芳香性。如果再引入其他原子，那么硼还能够形成更加丰富的多配位团簇构型，如图？所示，通过结构搜索算法研究人员找到了一系列金属-硼团簇。其中，Li2B24[]具有管状结构，具有三重芳香性，能够用于生长硼纳米管。碱土金属镁能够和硼形成具有高稳定性的鼓状结构MgB18团簇[]。铍元素能够和硼形成形似的团簇如具有二次旋转轴的BeB16-团簇[]。与硼同处于第13族的铝与硼有相同的价电子结构，铝元素对硼团簇结构的影响较小，但能够增加其稳定性，铝硼能够形成环状平面结构AlB18-[]，它由很大的带隙高的芳香性。以及一大类通过结构搜索方法找到的过渡金属-硼团簇[]。

在二维硼平面的结构搜索中，Wu等[]，使用第一性原理粒子群算法，预测了两种稳定的单层硼平面alpha1平面和beta1平面，他们的能量均比原先报道的由Tang等[]报道的alpha硼平面以及由Yakobson等[]报道发现的beta硼平面的结合能更低。利用全局结构搜索软件，Qu等[]预测发现稳定的TiB4结构。Zhang等[]报道发现了能带结构中具有狄拉克锥的FeB6结构。

可以看到，硼元素因其独特的缺电子特性，在结构多样性上异常丰富，但同样可以看到，由硼所形成的结构都是由硼元素以密堆的方式形成网络，但由于其缺电性，网络中需要空位或者其他元素原子作为电子给体来平衡电荷分布。根据以上特征，可以有选择的给出可能稳定的初始结构，使得所需要计算的结构数量大大降低而不再依赖于全局搜索算法。利用这样的方法，我们课题组在对平面团簇结构进行了高通量结构搜索，在数千个个候选结构进行能量的考察，不但能够完全找到已经被实验发现的平面团簇，还预测了稳定的B76团簇[]。对硼平面团簇使用类似的方法，找到了具有半导体性质的硼平面结构[]。

\section{论文的研究意义和结构安排}
二维硼材料在结构和性质上展现了丰富的特性，实验上研究人员已经能够生长出一定尺寸的单层硼材料，理论预测二维硼材料以及硼与其他元素所形成的二维化合物会表现出许许多多丰富的物理化学性质，尤其是和过渡金属原子形成的二维化合物理论计算预测力很多新奇的性质，使二维硼化合物不但有应用上的潜力，同时也是很好的新物理现象的研究对象。但硼在结构上有丰富的多样性，这一方面是其丰富的性质多样性的来源，但这也使得其合成和稳定结构的确定变得困难。如何发现和理论预测满足特定结构特定性质的二维硼材料，是本文研究的出发点。如果能从理论上预测具有优秀物理化学性质的二维硼材料，并对预测的材料的结构稳定性、结构特征、材料性质进行深入的研究，则能够快速地指导实验上此类材料的发现和应用，从而大大减少材料科学中实验的人工成本和材料成本。

针对二维硼材料结构上的共性，我们认为使用加入结构相关的先验知识进行有偏向性的筛选和搜索(biased-screening)的效果会优于使用全局搜索算法的无偏向性的结构搜索（unbiased-screening）。通过加入硼二维材料中硼原子之间倾向于保持包含空位的硼平面三角网格的结构，我们使用结构查重软件优先构建了尽可能均匀分布于势能面空间中目标结构能量局域极小值点附近的初始构型。再对所有的生成的结构进行精确的第一性原理计算获得准确的能量和性质。通过使用这个方法，我们研究了钛硼单层的稳定构型，不但发现了各个不同配比下已经报道过的几个钛硼单层材料，还发现了能量更低且具有优秀超导性质的TiB7单层构型。需要强调的是，通过加入先验的结构信息，我们所生成和计算的结构总数相比全局搜索算法所涉及的结构数量大大减少，该方法还更容易进行计算的管理，并且还获得了更好的结果。因此，我们认为对于特定材料，如果其结构随上拥有固定特征，可以考虑使用该方法首先生成合理构型，来快速的筛选符合指定目标的结构。

本论文将首先在第二章中介绍在论文相关的研究中所用到的理论方法，包括计算材料电子结构DFT理论和用于计算电声子互相作用的DFTP理论，以及通过第一性原理计算超导性质作用到的理论方法。第三章将重点阐述高通量材料计算的发展并引申出自动化工作流和数据库技术在高通量计算材料学中的重要性，最后集中描述AiiDA这一本论文研究中所用到的工作流工具的设计细节和使用方法。第四章描述如何对固定晶格的材料生成不重复的晶格构型所涉及的算法和软件实现。最后一章，则描述对钛硼构成的单层材料，利用第四章描述的工具进行不重复初始构型的生成，结构优化后得到局域稳定的构型，通过对结构的能量分析筛选可能稳定存在的构型并对其进行超导性质的研究。由于生成结构的数量非常巨大，在研究中我们使用了AiiDA作为高通量工具来自动化管理和生成计算得到的数据和计算过程数据库。

%%%%%%%%%%%%%%%%%%%%%%%%%%%%%%%%%%%%%%%%%%%%%%%%%%%%%%%%%%%%%%%%%%%%%%%%%%%%%%%
\chapter{电子结构理论与超导性质理论模拟}\label{chapter_theory}
% Define abbreviation
\newcommand{\ubm}[1]{\underline{\bm{#1}}}

在这一章中，我们将简要介绍凝聚态材料电子结构第一性原理计算所用到的理论方法。重点讨论
密度泛函理论（DFT）在材料基态性质以及密度泛函微扰理论（DFPT）在超导性质研究中的应用。
理论原理部分的讨论难以详尽，因此将仅涉及与论文内容相关的概念和讨论。

\section{密度泛函理论（DFT）及部分实现细节}

到目前为止，密度泛函理论（DFT）仍然是用于研究由电子和原子核通过库伦相互作用构成的多体系统的基态性质的最广泛实用的理论方法。该理论由Hohenberg 和 Kohn[] 于1964年以及Kohn 和 Sham[] 于1965年的两篇奠基性工作发展而来。简单来说，根据密度泛函理论，一个系统不但可以通过其多体波函数来完全描述，还可以通过其基态的电子密度来完全描述。体系中电子复杂的多体相互作用中难以处理的部分之间通过对总能的交换相关项贡献进行描述。这一交换相关项的具体形式当前并未知晓，但可以通过如基于自由电子气的模型近似来得到，常见的有局域密度近似（LDA）和广义梯度近似（GGA）。事实上，归功于计算能力的发展，DFT在近几十年中成为真实材料基态性质定量预测的一个极为精确有效的理论工具。

Hohenberg-Kohn（HK）定理描述了这样一个事实，即体系的基态电子密度$n(\bm{r})$与外势$V\mathrm{ext}$ 之间存在一一对应的关系。体系的电子密度由其所处的外势唯一确定，同时该电子密度也仅对应于该外势。因此，体系的能量是电子密度的泛函，该泛函的极小值对应于基态电子密度且给出了正确的基态能量。然而，多体系统的准确泛函形式并不知晓，密度泛函理论的数值方法实现基于Kohn-Sham（KS）方法。KS方法使用一个假想的辅助独立无粒子间相互作用的系统，以构成与真实多体系统相同的电子密度。从而，将一个有粒子间相互作用的多体问题转化为一组独立粒子与外场势之间可互相自洽求解的问题（类似平均场方法）。求解体系总能对电子密度泛函的最小极值，可以得到一系列单电子薛定谔函数，也就是KS方程。多体系统的电子密度因此可以通过求解KS方程得到:
\begin{equation}
  \left\{ -\nabla^2 + v_{\mathrm{eff}}(\bm{r}) \right\}\psi_i(\bm{r}) = \epsilon_i \psi_i (\bm{r}).
\end{equation}
$\epsilon_i$是方程的本征值，$\psi_i(\bm{r})$是对应的本征函数，有效势$v_{\mathrm{eff}}$ 是电子密度的泛函由外势能和诱导的屏蔽势能组成：
\begin{equation}
  v_{\mathrm{eff}}[n] = v_{\mathrm{ext}} + v_{\mathrm{scr}}[n] =
  v_{\mathrm{ext}} + v_{H}[n] + v_{XC}[n].
\end{equation}
电子密度$n$为:
\begin{equation}
  n(\bm{r}) = \sum_i f_i |\psi_i(\bm{r})|^2
\end{equation}
其中$f_i$为波函数$\psi_i$的占据函数。KS方程需要通过自洽的方式求解，因为KS哈密顿量$H^{KS}$中的有效Kohn-Sham势部分$v_{\mathrm{eff}}[n]=v_{\mathrm{ext}} + v_{H}[n] + v_{XC}[n]$是由电子密度确定的。$v_{\mathrm{ext}}$为电子受到的与原子核间的库伦相互作用，$v_{H}[n]$为电子间的Hatree势，而$v_{XC}[n]$则描述未知的交换-关联势，被定义为交换-关联能对电子密度的泛函微分。在局域密度近似（LDA）中，交换-关联能被认为仅依赖于局域的电子密度，且与均匀的自由电子气模型的交换关联能函数相等。而在广义梯度近似（GGA）中，交换-关联能还依赖于电子密度空间中的梯度。

我们可以看到，KS方法构建了与真实系统相同的基态电子密度和基态总能，但无法定量计算其他的一些性质，比如，KS波函数通常不能够准确对应光谱研究中的准粒子态，又或对于晶格变化导致的外场势变化引起的准粒子态无法单纯通过DFT理论进行描述。

周期性体系数值求解KS方程的常用方法是基于平面波的赝势近似方法。该论文的所有工作均用到这一方法，因此我们将对其进行简要的介绍。

\subsection{赝势近似}
在赝势近似中，我们将核与内层束缚核电子的库伦势作替换将势替换为仅作用于外层价电子的有效势$v^\mathrm{PS}$。通过近核处平滑的赝波函数，避免了需要过多平面波以描述剧烈振荡的近原子核处真实电子波函数的问题节约了计算成本。赝势近似前提假定核电子对化学键及化学特性没有重要的影响，且其并不随结构的变化而发生变化，因此可合理地认为核电子是牢牢固定并静止（frozen）与原子核附近，即静止核近似（frozen core approximation）。这一近似大大节约了计算成本，但代价是赝势在空间中是非局域的，
\begin{equation}
  v^{\mathrm{PS}}(\bm{r},\bm{r'}) = v^{\mathrm{loc}}(\bm{r}) + v^{\mathrm{nl}}(\bm{r},\bm{r'}).
\end{equation}
公式中$v\mathrm{loc}$表示局域贡献的部分，$v^\mathrm{nl}$代表非局域的部分。
第一性原理的赝势来自于原子全电子波函数求解，生成赝势的目标为构建能够几乎完全复现真实原子势散射特性的赝势。也就是超出特定原子核距离，将真实原子势替换为构建的赝势，有完全相同的原子能级本征值。

赝势的种类主要有以下几类：模守恒赝势（NC）[]，超软赝势（US）[]，的projecter-augmented wave（PAW）[]。 在NC赝势的构建中，要求赝波函数的模与对应的全电子波函数完全相等。这一条件在US和PAW的构建中适当放宽，从而可以产生更加平滑的赝波函数，但代价是US赝势和PAW赝势的数学表达更加复杂。

\subsection{平面波基组}
在处理周期性体系时，Kohn-Sham方程满足布洛赫定理，因此波函数可以写为$\psi_{n\bm{k}}(\bm{r})=e^{i\bm{k}\cdot \bm{r}} u_{n\bm{k}}(\bm{r})$的形式，其中$n$和$\bm{k}$分别为能带的索引和晶格波矢。函数$u_{n\bm{k}}(\bm{r})$是以晶格为周期的周期函数，满足$u_{n\bm{k}}(\bm{r}+\bm{R})=u_{n\bm{k}}(\bm{r})$，R为晶体的布拉维格子。

在周期性体系中，赝势方法给出的KS方程的数值解波函数可以通过平面波基组展开：
\begin{equation}
  \psi_{n\bm{k}}(\bm{r}) = \sum_{\bm{G}}e^{i(\bm{k}+\bm{G})\cdot \bm{r}}c_{n\bm{k}}(\bm{G}).
\end{equation}
其中求和项$G$为倒空间中全部的倒格矢，$c_{n\bm{k}}(\bm{G})$则是对波函数$\psi_{n\bm{k}}(\bm{r})$进行该展开的傅立叶系数。在真实的数值模拟中，我们仅能展开为有限数量平面波，因此对波矢作以下截断：
\begin{equation}
  \frac{|\bm{k} + \bm{G}|^2}{2} < E_{\mathrm{cutoff}}.
\end{equation}
$E_{\mathrm{cutoff}}$是波函数的截断能，该参数控制了展开平面波的数量，也因此影响了计算的精度。
波函数写为平面波基组展开的形式，KS方程$H\left|\psi_i\right>=\epsilon_i\left|\psi_i\right>$可写为矩阵形式的久期方程，本征值和本征波函数可通过迭代对角化得到。

\section{密度泛函微扰理论（DFPT）}

电子和声子可以被认为是固体的最基本的构筑单元，任何固体物质都包含了这两种元激发。对固态物质性质的本质的理解都基于对电子和声子这两种准粒子的清晰认识上。然而，要解出周期性体系中电子和原子核（或将原子核和内层非价电子构成的离子看作独立的对象）互相耦合的多体量子力学问题是一个接近不可能的任务。但由于电子和原子核在质量上的巨大差别，它们两者可以近似地看作独立的两个动力学子系统而独立求解。近几十年，快速发展的数值方法，使得从第一性原理求解电子部分的量子问题成为可能。大多数的第一性原理求解电子量子问题的方法都是基于上一节所描述的密度泛函的方法，现在有许多DFT的软件，使得电子结构的求解成为物质研究的一个例行的探索步骤。对声子部分的第一性原理求解的方法的成熟比电子部分的更晚，这是显而易见的，因为对基础的振动性质的精确的数值描述是要以精确的电子结构的描述为基础的。30年前所发展的通过微扰的线性响应的方法，也就是题目中的密度泛函微扰理论（DFTP）现在已经成为准确地描述声子的高效的数值方法。

电子和声子之间的相互作用即电声耦合，影响甚至主导了许多固相物质的物理性质。该相互作用最明显的例子是，在金属中，晶格振动影响了低能级的电子激发，从而导致了可观测的电子迁移性质和热力学性质的改变。而且，电声耦合是电子自发成对的原因之一，这导致了超导这一宏观的物理现象。

这一节我们将介绍DFT框架下，第一性原理数值求解声子的电声耦合的方法。再过渡到声子调制的电子成对相互作用从而得以描述超导性质。在开始，有必要先简要叙述并给出一些数学描述以方便后续的介绍。这部分的剩余内容将首先给出电声哈密顿量的绝热近似的定义，再介绍晶格振动的唯像的描述。

我们考虑固体由电子和离子组成，此处的离子是原子核和非价层的核电子组成的统一整体。晶体中电子和离子的动力学可以由以下哈密顿量描述：
\begin{equation}
  \mathcal{H} = T_e + V_{ee} + T_i + V_{ii} + H_{e-i},
\end{equation}
其中$T_e$和$T_i$分别是电子和离子的动能项，$ V_{ee}$，$V_{ii}$分别表示电子和离子之间的库伦相互作用势，$H_{e-i}$则概括了电子和离子的相互作用。

由于电子的质量远远小于离子的质量，我们可以认为较轻的电子比更重的离子以快的多的速度移动，从而近似对耦合的多体薛定谔方程$\mathcal{H} \Psi(\ubm{r},\ubm{R})=\varepsilon \Psi(\ubm{r},\ubm{R})$ (下划线表示是一组粒子的坐标)进行简化而使得方程可以数值解。我们认为在离子发生振动产生移动后，电子几乎瞬间发生了重新排布。这样的现象可以通过引入微扰参数$\kappa$来描述离子产生的位移，在离子质量趋向于无穷大$M\rightarrow+\infty$时，该参数趋向于$\kappa\rightarrow 0$，这个描述最早由Born和Oppenheimer[]用于分子的振动，并由Chester和Houghton[]拓展到对固体（周期性）的描述。

设离子振动偏移平衡位置:
\begin{equation}
  \bm{R}_i = \bm{R}^0_i + \kappa \bm{u}_i.
\end{equation}
动能项中参数$\kappa$指数项为$-2$,势能项该参数指数项为$2$，令动能项和势能项该参数的阶相等，则得到$\kappa$的形式为：$\kappa=（）^{1/4}$，该参数除了对氢和氦两种元素，值均小于$0.1$。使用这个参数，可以用微扰法对哈密顿量和波函数进行系统的展开。最低阶项展开式中波函数可以写为$\Psi(\ubm{r},\ubm{R})=\chi(\ubm{R})\psi(\ubm{r};\ubm{R})$，可以看到电子波函数仅参数地依赖于原子坐标。因此电子波函数满足方程：
\begin{equation}\label{eq:dftp_ele_eigen}
  \left[ T_e + V_{ee} + H_{e-i}(\ubm{R}) \right]\psi_n(\ubm{r};\ubm{R}) = E_n(\ubm{R})\psi_n(\ubm{R}); \ubm{R}),
\end{equation}
对离子位置R的依赖统一通过He-i描述，离子的位置是下列方程的解：
\begin{equation}
  \left[ T_i + V_{ii} + E_n(\ubm{R}) \right] \chi(\ubm{R}) = \xi \chi(\ubm{R}).
\end{equation}
这个近似被成为绝热近似或波恩-奥本海默近似。这个近似将电子和离子进行了动力学的解耦，忽略了离子移动导致的电子的激发。在上式？中，En是电子系统的某个本征值，通常可认为是解耦后电子薛定谔方程式？的基态解，而且通常来说，有限温度的电子激发不会对离子移动造成什么影响。这一项包含了重要的价电子对离子实移动的屏蔽效应，虽然该效应对基态和对激发态是一样的。

为得到比以上绝热近似更精确的解，我们使用式(\ref{eq:dftp_ele_eigen})的波函数解为基组来展开原系统的波函数：
\begin{equation}
  \Psi_m(\ubm{r};\ubm{R}) = \sum_n \chi_{mn}(\ubm{R})\psi_n(\ubm{r};\ubm{R}).
\end{equation}
代入系统薛定谔方程Hphi=Ephi整理得到离子部分满足方程：
\begin{equation}
  \left[ T_i + V_{ii} + E_n(\ubm{R}) \right] \chi_{mn}(\ubm{R}) + \sum_{n'}\Delta H_{nn'}\chi_{mn'}(\ubm{R}) = \varepsilon_m \chi_{mn} (\ubm{R}).
\end{equation}
与绝热近似时离子位置满足的方程？相比多出了左边第二项，其中$\Delta H=\Delta H^{(1)} + \Delta H^{(2)}$：
\begin{equation}
  \Delta H_{nn'}^{(1)} = -\frac{1}{M}\sum_i \int dr^{3N}\psi_n^*(\ubm{r};\ubm{R})\nabla_{R_i}\psi_{n'}(\ubm{r};\ubm{R})\cdot \nabla_{R_i},
\end{equation}
\begin{equation}
  \Delta H_{nn'}^{(2)} = -\frac{1}{2M}\sum_i \int dr^{3N}\psi_n^*(\ubm{r};\ubm{R})\nabla^2_{R_i}\psi_{n'}(\ubm{r};\ubm{R}).
\end{equation}
这两项中都有电子波函数对离子实位置的偏导，因此包含了离子实位置移动所导致的电子子系统的激发情况。第一项展开的系数为$\kappa$，第二项为$\kappa^2$，同时由于$\kappa$值很小（除原子质量最小的几个原子需要特别考虑），因此第一项的值远大于第二项，主导了非绝热过程。非绝热修正项对结果波函数的影响为$\kappa^3$，对能量的影响为$\kappa^6$。$\kappa$的值仅依赖于质量比而与电子离子相互作用的强度无关。综上绝热近似无论对于近自由电子的系统还是价电子束缚在离子实周围的体系都适用。

接下来我们来看看晶格振动的唯象描述。在绝热近似下，我们通过以下有效势来给出离子的静力学和动力学结果：
\begin{equation}\label{eq:effective_potential}
  \Omega(\ubm{R}) = V_{ii}(\ubm{R}) + E_0(\ubm{R}),
\end{equation}
其中$E_0(\ubm{R})$是给定的一组离子坐标$\ubm{R}$下的电子基态能量。许多综述文章[][]都是从这个有效势$\Omega$开始来构建晶格动力学的微观理论。

我们通过对势能$\Omega$在参考位置$\ubm{R}^0$附近位移$u$来进行展开以得到系统动力学性质：
\begin{equation}
  \Omega(\ubm{R}) = \Omega(\ubm{R}^0) + \sum_{i\alpha} \Phi_\alpha(i)u_{i\alpha} + \frac{1}{2} \sum_{i\alpha j\beta} \Phi_{\alpha\beta}(i,j) u_{i\alpha}u_{j\beta} + \cdots .
\end{equation}
希腊字母$\alpha$和$\beta$描述的是离子的空间笛卡尔坐标序号即$x,y,z$，$i$和$j$是原子的序号。第一阶的系数$\Phi_\alpha$是作用在参考原子上的力的相反数：
\begin{equation}
  F_{i\alpha} = -\frac{\partial\Omega}{\partial R_{i\alpha}}\bigg|_0 = -\Phi_\alpha(i).
\end{equation}
如果我们选择参考位置为平衡位置，则此时该项为零，且平衡位置处势能$\Omega$达到极小值。二阶项系数为：
\begin{equation}
  \Phi_{\alpha\beta} = \frac{\partial^2\Omega}{\partial R_{i\alpha}\partial R_{j\beta}} \bigg|_0 .
\end{equation}
这一系数的物理意义为，当第$i$个原子偏离平衡位置所导致的第$j$个原子感受到的力的变化：
\begin{equation}
  F_{j\beta} = -\sum_{i\alpha} \Phi_{\alpha\beta}(i,j)u_{i\alpha}.
\end{equation}
因此该项描述了位移随力的线性变化。也就是弹簧振子的简谐振动力常数的三维（张量）版本，而更高次项系数被称为非简谐力常数。如果仅截断到二次项，就是我们常说的简谐近似。
在周期性体系中，原子的序号可以使用一对符号来表示即；$i=(l\kappa)$，其中$l$表示单胞序号，$\kappa$表示原子在单胞内的序号。通过傅立叶变化将力常数矩阵变换到倒空间为动力学矩阵：
\begin{equation}
  D_{\kappa\alpha\kappa'\beta}(\bm{q}) = \frac{1}{\sqrt{M_\kappa M_{\kappa'}}} \sum_l \Phi_{\alpha\beta}(l\kappa,0\kappa')e^{-i\bm{q}(\bm{R}^0_{l\kappa}-\bm{R}^0_{0\kappa'})},
\end{equation}
该矩阵决定了晶体中的振动模式即所谓的声子，
\begin{equation}
  \sum_{\kappa'\beta}D_{\kappa\alpha\kappa'\beta}(\bm{q})\eta_{\kappa'\beta}(\bm{q}j) = \omega^2_{\bm{q}j}\eta_{\kappa\alpha}(\bm{q}j).
\end{equation}
其中$\omega^{\bm{q}j}$和$\eta_{\kappa\alpha}(\bm{q}j)$分别是确定波矢$\bm{q}$第分支$j$的频率和振动模式。通过这两个量可以建立原子位移和声子之间的联系，设声子的升降算符为$b$和$b^{\dagger}$，则：
\begin{equation}
  u_{l\kappa\alpha} = e^{i\bm{q}\bm{R}^0_{l\kappa}}\frac{1}{\sqrt{N_q}} \sum_{\bm{q}j} A^{\bm{q}j}_{\kappa\alpha}(b_{\bm{q}j} + b^{\dagger}_{-\bm{q}j}),
\end{equation}
其中：
\begin{equation}
  A^{\bm{q}j}_{\kappa\alpha} = \frac{\eta_{\kappa\alpha}(\bm{q}j)}{\sqrt{2M_\kappa \omega_{\bm{q}j}}}.
\end{equation}
要完全的描述简谐振动的谱性质，需要知道全布里渊区的振动模式信息，或者知道所有原子键之间的力常数。在金属材料中，由于电子对离子移动的屏蔽，实空间中晶格间的键的相互作用很小，可以在比较小的距离截断，因而使用力常数的描述更节约计算资源。

\subsection{第一性原理晶格动力学}
晶格动力学的首要目标就是计算决定离子动力学的量的数值。式（\ref{eq:effective_potential}）中有效势分解为的左边第一项是离子间的库伦相互作用，可以直接得到其对力常数的贡献。第二项综合了由电子贡献的与键和电子对离子屏蔽的相关物理量的信息，该项需要通过准确的电子结构来给出，可以通过密度泛函理论（DFT）对体系基态求解实现。

晶格势决定了晶格的动力学性质，而晶格势等于固定原子位置后构型的基态能量，该能量可以通过DFT求解。从第一性原理求解晶格动力学性质主要可以分为两类方法：1）直接求解力矩阵，2）线性响应法。直接法是通过计算理想晶格以及计算特定的原子移动后的构型能量来得到力矩阵。冷冻声子法（FP）是概念上最为简单且最早被使用直接求力矩阵的方法，它通过能量和位移的二次关系来提取振动模式及模式的频率。这个方法需要首先确定声子的本征函数从而直到如何移动原子，因此该方法需要提前直到晶体的对称性信息。但由于FP方法中对原子移动了有限的距离，因此它可以包含非简谐的效应，可以通过拟合能量和力的高阶关系用来计算更高阶的系数。通过移动个别原子后其他单胞内原子受力改变和移动原子距离的线性关系来计算晶格性质是更为高效的一个方法。另一个更有效的直接求力矩阵的方法是，通过海尔曼-费曼定理，即原子受力可以直接从移动原子后体系基态求解过程中得到（不需要作能量对位置偏导）。因此对应一个原子的移动只需要简单的计算就可以给出动力学矩阵的一行，合适地选取原子位移，就能够得到整个动力学矩阵。该方法并不需要提前得到晶体的对称性信息。和线性响应法相比，直接法的缺点是如果要求非零波矢的声子（除$\Gamma$点外）性质，需要构建超胞，如果要得到完整的声子谱，则超胞需要大到晶格单胞之间的相互作用可以忽略不计。线性响应法通过微扰的方法来计算能量的导数，力学矩阵可以从二阶偏导（公式？）直接得到。使用这个方法，无需构建超胞就能直接在倒空间任意波矢求解曾而获得完整的动力学性质。因为在本论文工作的稳定性确定和超导计算部分用到了这以方法，我们在下面详细描述该方法。

\subsection{线性响应方程}
本小节我们给出在DFT框架下周期体系基于对原子位置微扰的方法求动力学矩阵的详细过程。首先我们假设这样的情况，即外场势$v_{\mathrm{ext}}$依赖于一组绝热的微扰参数$\Lambda={\lambda_a,\lambda_b \cdots \lambda_p}$，每个外势$v_{\mathrm{ext}}^\Lambda$对应一个电子密度和确定的总能
\begin{equation}
  E^\Lambda = F[n^\Lambda]+\int d^3 r n^\Lambda (\bm{r})v_\mathrm{ext}^\Lambda (\bm{r}).
\end{equation}
因此加外势扰动后总能对微扰参数的偏导为：
\begin{equation}
  \frac{\partial E^\Lambda}{\partial \lambda_a} =
  \int d^3 r n^\Lambda(\bm{r})\frac{\partial v^\Lambda_\mathrm{ext}(\ubm{r})}{\partial \lambda_a} +
  \int d^3 r \frac{\delta E^\Lambda}{\delta n(\bm{r})} \frac{\partial n^\Lambda(\bm{r})}{\partial \lambda_a}.
\end{equation}
第二项的电子密度为势能量最小的基态电子密度，由变分法可知第二项为零。而这就是DFT框架下的海尔曼-费曼定理[]。微扰后能量的二阶偏导为：
\begin{equation}\label{eq:second_order_energy}
  \frac{\partial^2 E^\Lambda}{\partial\lambda_a\partial\lambda_b} =
  \int d^3 r\frac{\partial n^\Lambda(\bm{r})}{\partial\lambda_b}\frac{\partial v^\Lambda_\mathrm{ext}(\bm{r})}{\partial\lambda_a} +
  \int d^3 r n^\Lambda(\bm{r})\frac{\partial^2 v^\Lambda_\mathrm{ext}(\bm{r})}{\partial\lambda_a\partial\lambda_b}.
\end{equation}
从上式可知，要求能量的二阶导只需要得到关于电子密度的一阶导。因此仅需要考虑系统对微扰的线性响应。有效势在微扰后可以写成以下形式：
\begin{equation}
  \delta v_{\mathrm{eff}}(\bm{r}) = \delta v_\mathrm{ext}(\bm{r}) + \delta v_\mathrm{scr}(\bm{r}) =
  \delta v_\mathrm{ext}(\bm{r}) + \int d^3 r' I(\bm{r},\bm{r'})\delta n(\bm{r'})
\end{equation}
\begin{equation}
  I(\bm{r},\bm{r'}) \equiv \frac{\delta v_\mathrm{scr}(\bm{r})}{\delta n(\bm{r'})} = \frac{\delta v_H(\bm{r})}{\delta n(\bm{r'})} +
  \frac{\delta v_{XC}(\bm{r})}{\delta n(\bm{r'})} =
  \frac{2}{|r-r'|} + \frac{\delta^2 E_{XC}}{\delta n(\bm{r}) \delta n(\bm{r'})}.
\end{equation}
使用微扰原理，KS波函数的一阶变分为：
\begin{equation}
  \delta \psi_i(\bm{r}) = \sum_{\substack{i,j\\i\neq j}} \frac{\langle i |\delta v_{\mathrm{eff}}|j \rangle}{\epsilon_i-\epsilon_j} \psi_j(\bm{r})
\end{equation}
同样写出$\delta \psi_i^*(\bm{r})$可以得到电子密度的一阶变分为：
\begin{align}
  \delta n(\bm{r}) &= \sum_i f_i [\psi_i^*(\bm{r})\delta \psi_i(\bm{r}) + \delta \psi_i^*(\bm{r}) \psi_i(\bm{r})] \\
  &= \sum_{i\neq j} \frac{f_i-f_j}{\epsilon_i-\epsilon_j} \langle j| \delta v_\mathrm{eff} | i \rangle \psi_i^*(\bm{r}) \psi_j(\bm{r}).
\end{align}

从公式？和？可以发现，$\delta n$和$\delta_\mathrm{eff}$互相依赖，需要通过自洽法求解。接着为了凸显电子密度随有效势$\delta_\mathrm{eff}$的变化关系我们还可以把以上关系写为：
\begin{equation}
  \delta n(\bm{r}) = \int d^3 r' \chi_0 (\bm{r},\bm{r'})\delta v_\mathrm{eff}(\bm{r'})
\end{equation}
\begin{equation}
  \chi_0(\bm{r},\bm{r'}) = \sum_{i\neq j}\frac{f_i-f_j}{\epsilon_i-\epsilon_j} \psi^*_i(\bm{r})\psi_j(\bm{r})\psi^*_j(\bm{r'})\psi_i(\bm{r'}).
\end{equation}
此处的$\chi_0$是无相互作用的KS系统的电极化率，它仅仅依赖于系统基态。虽然使用微扰方法，但由于处理的是电子无相互作用的KS系统，因此没有更高阶的近似。
现在已知有效势对电子密度的影响，要得到电子密度和外势之间的关系，我们使用式(?)有效势可以写为：
\begin{align}
  \delta v_\mathrm{eff} &= \delta v_\mathrm{ext} + I \chi_0 \delta v_\mathrm{eff} \\
  &= [1-I\chi_0]^{-1} \delta v_\mathrm{ext} = \epsilon^{-1} \delta v_\mathrm{ext}.
\end{align}
其中$\epsilon$为静态介电常数矩阵，描述的是体系对外场微扰的屏蔽。

为高效求解该介电常数矩阵的逆就引出了本章所要介绍的方法DFTP，其克服了许多其他方法的不足，如求解$\chi_0$时需要未占据态，以及对于周期体系实空间中的介电矩阵为无穷维，通过傅立叶变换到倒空间该矩阵逆的求解十分耗费计算资源。DFTP通过自洽法来实现线性响应法而避免了介电常数矩阵求解时的困难。

公式030电子密度的一阶变分中包含了对电子态的两次的求和，而参数$(f_i-f_j)/(\epsilon_i-\epsilon_j)$限制了（对于半导体，对金属的拓展可参考文献[]）其中一个态为导带另一个态来自价带，不考虑磁性时体系满足时间反演对称性，式030可以写为：
\begin{equation}
  \delta n(\bm{r}) = 2\sum_{vc}\frac{1}{\varepsilon_v-\varepsilon_c}
  \langle c| \delta v_\mathrm{eff} | v\rangle \psi^*_v(\bm{r})\psi_c(\bm{r}).
\end{equation}
整理并把导带部分写在一起后得：
\begin{equation}
  \delta n(\bm{r}) = 2\sum_v \psi^*_v(\bm{r}) \Delta_v(\bm{r})
\end{equation}
\begin{equation}
  |\Delta_v \rangle = \sum_{c}\frac{1}{\varepsilon_v-\varepsilon_c}
  |c\rangle \langle c| \delta v_\mathrm{eff} | v\rangle
\end{equation}
该右矢$|\Delta_v\rangle$满足下列线性方程：
\begin{equation}
  (H-\epsilon_v)|\Delta_v\rangle = -\sum_c |c\rangle \langle c| \delta v_\mathrm{eff} | v\rangle = (P_v-1)\delta_\mathrm{eff}|v\rangle .
\end{equation}
因此，给定$\delta v_\mathrm{eff}$可以从上述线性方程组求得$|\Delta_v\rangle$，代回031可以得到$\delta n$，进而能得到新的$\delta v_\mathrm{eff}$。自洽得到最后的结果。用该自洽方法求解外场对电子密度的扰动{\textbf{不需要}}用到导带部分的信息。

\subsection{DFPT计算声子性质}
这里我们用DFPT来计算原子间力常数。将全实空间的偏离平衡位置的位移写成下列函数的形式，从而可以描述特定声子对应实空间的不同单胞的原子的位移：
\begin{equation}
  u_{l\kappa\alpha} = d_{\kappa\alpha}e^{i\bm{q}\bm{R}^0_{l\kappa}} + d^*_{\kappa\alpha}e^{-i\bm{q}\bm{R}^0_{l\kappa}},
\end{equation}
其中$l$是单胞的序号，$\kappa$是单胞内的某个原子，$\alpha$是笛卡尔坐标序号$x$，$y$或$z$。复数表示的振幅$d_{\kappa\alpha}$使得可以包含特定的原子位移的相对相位。动力学矩阵可以看作是势能面能量对特定波矢的声子的振幅$$d_{\kappa\alpha}$$的二阶偏导：
\begin{equation}
  D_{\kappa\alpha\kappa'\beta}(\bm{q}) =
  \frac{1}{\sqrt{M_\kappa M_{\kappa'}}} \delta^{\bm{q}}_{\kappa\alpha} \delta^{-\bm{q}}_{\kappa'\beta} E \bigg|_{\bm{u}=0}.
\end{equation}
其中$\delta^{\bm{q}_{\kappa\alpha}}\equiv \frac{\partial}{\partial d_{\kappa\alpha}}$以及$\delta^{\bm{-q}_{\kappa'\beta}}\equiv \frac{\partial}{\partial d^*_{\kappa\alpha}}$。
使用式（\ref{eq:second_order_energy}），动力学矩阵可以分为两部分，即来自于离子势的二阶偏导和电子的一阶偏导所贡献的屏蔽部分。
通常在晶体中，外势可以写为由原子势叠加而成的形式：
\begin{equation}
  v_\mathrm{ext}(\bm{r}) = \sum_{l\kappa} v_\kappa (\bm{r}-\bm{R}_{l\kappa}).
\end{equation}
平衡位置处一阶偏导的值为：
\begin{align}
  \delta^{\bm{q}}_{\kappa\alpha}v_{\mathrm{ext}}(\bm{r}) &=
  - \sum_l \nabla^{\bm{r}}_\alpha v_\kappa (\bm{r}-\bm{R}^0_{l\kappa})e^{i\bm{q}\bm{R}^0_{l\kappa}} \\
  &= -e^{i\bm{qr}} \sum_l e^{i\bm{q} (\bm{R}^0_{l\kappa}-\bm{r})} \nabla^{\bm{r}}_\alpha v_\kappa(\bm{r}-\bm{R}^0_{l\kappa}).
\end{align}
而电子部分的一阶导，则就需要用到上一小节介绍的DFPT方法。首先将电子密度的变分通过傅里叶变换到倒空间为：
\begin{equation}
  \delta^{\bm{q}}_{\kappa\alpha} n(\bm{q}+\bm{G}) = - \frac{4}{V}\sum_{\bm{k}v} \langle \bm{k}v| e^{-i(\bm{q}+\bm{G})\bm{r}}| \Delta^{\bm{q}}_{\kappa\alpha}(\bm{k}v) \rangle,
\end{equation}
$V$为晶体的单胞体积，其中的Delta即为式031中的部分，定义如下：
\begin{equation}
  | \Delta^{\bm{q}}_{\kappa\alpha}(\bm{k}v) \rangle =
  \sum_c \frac{|\bm{k}+\bm{q}c\rangle \langle \bm{k}+\bm{q}c | \delta^{\bm{q}}_{\kappa\alpha} v_\mathrm{eff} | \bm{k}v \rangle}
  {\epsilon_c(\bm{k}+\bm{q})-\epsilon_v(\bm{k})}.
\end{equation}
可以通过求解下列非齐次线性方程组得到：
\begin{equation}
  (H^{\bm{k}+\bm{q}}_{KS}-\epsilon_v(\bm{k}))|\Delta^{\bm{q}}_{\kappa\alpha}(\bm{k}v)\rangle =
  (P^{\bm{k}+\bm{q}}_v-1)\delta^{\bm{q}}_{\kappa\alpha}v_\mathrm{eff}|\bm{k}v\rangle .
\end{equation}
式041,042,043组成的三个方程可以自洽求解。自洽后得到特定倒空间点$\bm{q}$的电子和外式的变分，从而可代入求解该$\bm{q}$点处的动力学矩阵得对应的声子模式。
\begin{equation}
  \delta^{\bm{q}}_{\kappa\alpha}\delta^{-\bm{q}}_{\kappa'\beta} E =
  \sum_{\bm{G}} [ \delta^{\bm{q}}_{\kappa\alpha}n(\bm{G}+\bm{q})\delta^{\bm{-q}}_{\kappa'\beta}v_\mathrm{ext}(\bm{G}+\bm{q}) + \delta^{\bm{q}}_{\kappa\alpha}\delta^{\bm{-q}}_{\kappa'\beta} v_\mathrm{ext}(\bm{G})].
\end{equation}

\section{电声耦合和超导性质的计算}
从Frohlich哈密顿量[]出发我们可以得到电子之间的相互作用满足以下关系：
\begin{equation}
  V_{\mathrm{eff}}(\bm{k},-\bm{k},\bm{q}) = |g_{\bm{k},\bm{q}}|^2
  \frac{\omega_{\bm{q}}}{(\epsilon_{\bm{k}}-\epsilon_{\bm{k}+\bm{q}})^2-\omega^2_{\bm{q}}}
\end{equation}
其中$g_{\bm{k},\bm{q}}$为耦合函数描述了体系中电子和声子的耦合，可以通过下面关系得到，所需要的参数都可以从体系电子结构和声子谱的第一性原理计算中得到：
\begin{equation}\label{eq:coupling_coeff}
  g^{\bm{q}\lambda}_{\bm{k}+\bm{q}\nu',\bm{k}\nu} =
  \sum_{\kappa\alpha} A^{\bm{q}j}_{\kappa\alpha}
  \langle {\bm{k}+\bm{q}\nu'} | {\delta^{\bm{q}}_{\kappa\alpha}v_{\mathrm{eff}}} | {\bm{k}\nu} \rangle .
\end{equation}

超导性质是电子系统体现的宏观现象。它源于由于低于临界温度时，体系所处的费米液体态不稳定而向更稳定的新的电子和电子成对（库伯对）的基态过渡。BCS理论[]展示的是，当体系中两个电子存在吸引相互作用时，该电子电子成对的状态是稳定的。而这个电子电子成对的吸引力就是由电声耦合提供的。电声耦合（EPC）是已知的大多数的超导现象的来源，这类超导通常成为常规超导体以区别其他材料中存在的尚未有完全清楚的理论机制的成对相互作用而导致的超导。

BCS理论，假设了电声耦合是弱耦合。后续在BCS基础上应用多体系统处理的方法[]发展而来的Elishberg理论[]将BCS理论中的弱耦合拓展到了强耦合，从而可以定量的预测许多超导态的性质。超导态的一个重要性质是其准粒子的谱上有一个带隙，带隙的大小是超导性质的一个重要的序参量。Elishberg理论的一个重要的结果是其描述超导能带的方程的计算中中不需要用到体系的超导态而只需要特定体系的正常电子态的信息。包括体系额的电子结构，的声子调制的电子成对能，这些信息通过第一性原理和上面的电声耦合的计算就能得到。

在本论文中我将省略Eliashberg理论的导出而仅仅列出其在超导性质计算中所用到的结果，和有效电子电子相互作用的计算过程。声子调制的电子电子相互作用是两个电子之间通过交换一个声子的物理过程，可以由下列的Eliashberg函数表示：
\begin{equation}
  \alpha^2 F_{\bm{k}\nu,\bm{k}'\nu'}(\omega) =
  N(\epsilon_F) \frac{1}{N_{\bm{q}}} \sum_{\bm{q}j}
  |g^{\bm{q}j}_{\bm{k}'\nu',\bm{k}\nu}|^2 \delta(\omega-\omega_{\bm{q}j}).
\end{equation}
其中的$N(\epsilon_F)$是费米能级处单个电子的电子态密度。在固定频率$\omega$对所有该过程涉及的声子模式求和。耦合的相互作用就是前面提到的包含屏蔽效应的电声相互作用耦合函数(矩阵元表示)即式（\ref{eq:coupling_coeff}）。该式右边只设计非超导态的性质。又由于动量守恒，在该过程中$\bm{q}=\bm{k}-\bm{k}'$。且该电子成对相互作用形成的过程过程中，有效的电子态集中在能量为，$|\epsilon_{\bm{k}\nu}-\epsilon_F|<\omega_{\mathrm{phonon}}$处，因此在式？中可以仅仅考虑费米面处的电子[]。

大多数超导体都表现出各向同性的超导带隙，原因是真实材料中都存在缺陷，而缺陷的存在会抵消电子相互作用的动量依赖。根据这两个假设，在大多数情况下可以近考虑费米面附近的电子态，且使用各项同性的Elishaberg函数：
\begin{equation}\label{eq:elishaberg_iso}
  \alpha^2 F(\omega) = \frac{1}{N(\epsilon_F)}\frac{1}{N_{\bm{q}}}
  \sum_{\substack{\bm{q}j \\ \bm{k}\nu\nu'}} |g^{\bm{q}j}_{\bm{k}+\bm{q}\nu',\bm{k}\nu}|^2
  \delta(\omega-\omega_{\bm{q}j})
  \delta(\epsilon_{\bm{k}\nu}-\epsilon_{F})
  \delta(\epsilon_{\bm{k}+\bm{q}\nu'}-\epsilon_F).
\end{equation}
在Elishaberg理论中重要的宏观物理量如超导转变温度都是对全频率的积分形式。各项同性的耦合常数可以表示为：
\begin{equation}\label{eq:lambda_coupling}
  \lambda = 2 \int d\omega \frac{\alpha^2 F(\omega)}{\omega}.
\end{equation}
该量是一个无量纲的量，描述了体系电声耦合的平均强度。由积分内的$1/\omega$可以看出在电子成对形成过程中低频的声子的贡献要比高频的声子的贡献更大。

利用$T\rightarrow 0$时的声子线宽$\gamma_{\bm{q}j}$的表达式（定义和推导见附录？），可以将Elishaberg函数写为下列我们在最后一章钛硼二维材料超导性质计算中用到的形式：
\begin{equation}\label{eq:elishaberg_calc}
  \alpha^2 F(\omega) = \frac{1}{2\pi N(\epsilon_F)}\frac{1}{N_{\bm{q}}}
  \sum_{\bm{q}j} \frac{\gamma_{\bm{q}j}}{\omega_{\bm{q}j}}
  \delta(\omega-\omega_{\bm{q}j}),
\end{equation}
代入式（\ref{eq:lambda_coupling}）得耦合常数：
\begin{equation}
  \lambda = \frac{1}{\pi N(\epsilon_F)}\frac{1}{N_q}
  \sum_{\bm{q}j}\frac{\gamma_{\bm{q}j}}{\omega^2_{\bm{q}j}}.
\end{equation}
从这个形式可以看出，Elishaberg函数是对所有的声子分支进行积分中再对声子的动量求平均。

在用DFPT方法计算Elishaberg函数的实际操作中，首先对声子线宽进行计算再对其在全部声子频率求和得到。式（\ref{eq:elishaberg_calc}）中的$\delta$函数在实际数值计算时通常使用平滑的高斯函数来代替，为了达到该近似的收敛往往需要很密的$k$点。而很密的$k$的计算可以使用Wannier函数来得到，该方法被称为电声耦合Wannier函数方法[]。

Frohlich哈密顿量的电子部分没有考虑电子电子本来的库伦相互作用（哈密顿量中电子部分由单粒子算符表示）。电子电子间的库伦相互作用在结果中主要影响的是式（\ref{eq:elishaberg_iso}）中准粒子的本征值$\epsilon$以及相应的声子的本征值$\omega$。实际物理过程包含有电子间的库伦相互作用，因此我们不能在声子调制的超导性质的计算中忽略这个相互作用，因为这个相互作用是排斥的因此会大大抑制前面导出的成对相互作用（吸引）。与电声耦合参数类似，可以引入所谓库伦参数来描述这个相互作用，即费米面处的平均有效库伦屏蔽。因为电子比晶格振动在时间尺度上快的多，因此能量上也比声子大的多，我们再定义Morel-Anderson库伦赝势[]来描述不同声子频率上限下的库伦参数：
\begin{equation}
  \mu = N(0)\langle\langle {V_C(k,k')} \rangle\rangle_{FS}
\end{equation}
\begin{equation}
  \mu^*(\omega_c) = \frac{\mu}{1+\mu \ln(\varepsilon_0/\omega_c)}.
\end{equation}
$\varepsilon_0$为电子系统的特征能量尺度，同时引入了一个比声子能大得多的截止$\omega_c$。在本论文的工作中我们使用了经验值来设置该参数$\mu^*$。实际上还可以通过更精确的超导第一性原理来计算得到这个参数[]。

超导转变温度$T_c$仅仅依赖于与材料相关的物理量$\alpha^2 F(\omega)$和$\mu^*$的值。Allen和Dynes[]对各向同性带隙方程进行了深入的数值分析，他们使用了一套$\alpha^2 F(\omega)$的标准谱，并在很大的参数范围内改变$\lambda$和$\mu^∗$。他们的研究揭示了两个重要结论，首先，他们发现在一个较小的参数空间(即$\lambda<2$和$\mu^∗<0.15$)中，$T_c$可以很好地近似于最初由McMillan[]提出的$T_c$公式，不同的是需要加入一个修正的前因子：
\begin{equation}
  T_c = \frac{\omega_{\mathrm{log}}}{1.2}
  \exp{\left[ {-\frac{1.04(1+\lambda)}{\lambda-\mu^*(1+0.62\lambda)}} \right]}
\end{equation}
前置因子包含了声子谱的平均频率（以对数方式定义的平均值），并与耦合强度加权平均，
\begin{equation}
  \omega_\mathrm{log} = \exp \left[ {\int d\omega \log(\omega)W(\omega)} \right],
\end{equation}
其中的权重$W$定义为：
\begin{equation}
  W(\omega) = \frac{2}{\omega} \frac{\alpha^2 F(\omega)}{\omega}.
\end{equation}
该Tc公式是对弱耦合极限下BCS公式$T_c = 1.13 \omega_D \exp(−1/\lambda)$的改进。

%%%%%%%%%%%%%%%%%%%%%%%%%%%%%%%%%%%%%%%%%%%%%%%%%%%%%%%%%%%%%%%%%%%%%%%%%%%%%%%
% 学位论文的正文应以《结论》作为最后一章
\chapter{结论}\label{chapter_concludes}

本文在第\ref{chapter_smallworld}章中，通过考虑数据中心网络布局构建中的最大度限制
问题，提出了符合数据中心网络基本要求的DS小世界模型，并分析了它的性质。随后提出
SIDN，将DS模型映射到具体的网络结构中，并分析了所构成网络的平均直径、网络总带宽、
对故障的容错能力等各项网络性能。

分析与仿真实验证明，SIDN网络具有很好的扩展能力，网络总带宽与网络规模成
近似线性增长的关系；具有很强的容错能力，链路损坏与节点损坏几乎无法破坏
网络的联通性，故障率对网络性能的影响与破坏节点/链路占总资源比率线性相关。

随后在第\ref{chapter_scalefree}章中，分析了无尺度网络在数据中心网络构建应用中的
理论方面问题。对Scafida \cite{gyarmati2010scafida}文中所述在最大度限制的情况下运
用BA算法构造的网络并不会损失无尺度性质的观点，进行了深入的分析，并指出了该论点的
局限性。

在给出了在引入节点最大度限制之后，利用分治和递归的思想，对无尺度网络
进行多层构建，对所构造的网络进行度-度相关性，以及聚类性分析。

\begin{table}
  \centering
  \begin{tabular}{cccp{38mm}}
    \toprule
    \textbf{文档域类型} & \textbf{Java类型} & \textbf{宽度(字节)} & \textbf{说明} \\
    \midrule
    BOOLEAN  & boolean &  1  & \\
    CHAR     & char    &  2  & UTF-16字符 \\
    BYTE     & byte    &  1  & 有符号8位整数 \\
    SHORT    & short   &  2  & 有符号16位整数 \\
    INT      & int     &  4  & 有符号32位整数 \\
    LONG     & long    &  8  & 有符号64位整数 \\
    STRING   & String  &  字符串长度  & 以UTF-8编码存储 \\
    DATE     & java.util.Date & 8 & 距离GMT时间1970年1月1日0点0分0秒的毫秒数 \\
    BYTE\_ARRAY & byte$[]$ & 数组长度 & 用于存储二进制值 \\
    BIG\_INTEGER & java.math.BigInteger & 和具体值有关 & 任意精度的长整数 \\
    BIG\_DECIMAL & java.math.BigDecimal & 和具体值有关 & 任意精度的十进制实数 \\
    \bottomrule
  \end{tabular}
  \caption{测试表格}\label{table:test5}
\end{table}

表\ref{table:test5}用于测试表格。随后分析了无尺度网络构造过程中，交换机节点与数
据节点的角色区别，分析了两者在不同比率下形成的网络形态，以及对网络性能造成的影响。

通过理论分析和仿真实验，分析并找出比率因子q的最佳取值。此外，无尺度现象
的引入提高了网络的聚类系数，从而在不失灵活性可靠性的基础上，进一步提升
了网络的性能。

在第\ref{chapter_random}章中，将关注点转移到交换机本身。由于图论难以描述数据中心
网络中的交换设备，因此放弃基于图的抽象模型，转而基于多维簇划分的思想，提出并设计
了WarpNet网络模型。

该网络模型突破了基于图描述的局限性，并对网络的带宽等指标进行理论分析并
给出定量描述。最后对比了理论分析、仿真测试结果，并在实际物理环境中进系
真实部署，通过6节点的小规模实验以及1000节点虚拟机的大规模实验，表明该模
型的理论分析、仿真测试与实际实验吻合，并在网络性能、容错能力、伸缩性灵
活性方面得到了进一步的提升。

在第\ref{chapter_experiments}章中，针对网络模型研究这一类工作的共性，设计构造通
用验证平台系统。以海量虚拟机和虚拟分布式交换机的形式，实现了基于少量物理节点，对
大规模节点的模拟。其模拟运行的过程与真实运行在实现层面完全一致，运行的结果与真实
环境线性相关。除为本文所涉若干网络模型提供验证外，可进一步推广到更为广泛的领域，
为各种网络模型及路由算法的研究工作，提供分析、指导与验证。

%%%%%%%%%%%%%%%%%%%%%%%%%%%%%%%%%%%%%%%%%%%%%%%%%%%%%%%%%%%%%%%%%%%%%%%%%%%%%%%
% 致谢，应放在《结论》之后
\begin{acknowledgement}
  首先感谢我的母亲韦春花对我的支持。其次感谢我的导师陈近南对我的精心指导和热心帮助。接下来，
  感谢我的师兄茅十八和风际中，他们阅读了我的论文草稿并提出了很有价值的修改建议。

  最后，感谢我亲爱的老婆们：双儿、苏荃、阿珂、沐剑屏、曾柔、建宁公主、方怡，感谢
  你们在生活上对我无微不至的关怀和照顾。我爱你们！
\end{acknowledgement}

%%%%%%%%%%%%%%%%%%%%%%%%%%%%%%%%%%%%%%%%%%%%%%%%%%%%%%%%%%%%%%%%%%%%%%%%%%%%%%%
% 附录
\appendix

\chapter{博士(硕士)学位论文编写格式规定(试行)}

\section{适用范围}

本规定适用于博士学位论文编写，硕士学位论文编写应参照执行。

\section{引用标准}

GB7713科学技术报告、学位论文和学术论文的编写格式。

GB7714文后参考文献著录规则。

\section{印制要求}

论文必须用白色纸印刷，并用A4(210mm×297mm)标准大小的白纸。纸的四周应留足空白
边缘，上方和左侧应空边25mm以上，下方和右侧应空边20mm以上。除前置部分外，其它
部分双面印刷。

论文装订不要用铁钉，以便长期存档和收藏。

论文封面与封底之间的中缝（书脊）必须有论文题目、作者和学校名。

\section{编写格式}

论文由前置部分、主体部分、附录部分(必要时)、结尾部分(必要时)组成。

前置部分包括封面，题名页，声明及说明，前言，摘要(中、英文)，关键词，目次页，
插图和附表清单(必要时)，符号、标志、缩略词、首字母缩写、单位、术语、名词解释
表(必要时)。

主体部分包括绪论(作为正文第一章)、正文、结论、致谢、参考文献表。

附录部分包括必要的各种附录。

结尾部分包括索引和封底。

\section{前置部分}

\subsection{封面（博士论文国图版用）}

封面是论文的外表面，提供应有的信息，并起保护作用。

封面上应包括下列内容：
\begin{enumerate}
\item 分类号  在左上角注明分类号，便于信息交换和处理。一般应注明《中国图书资
  料分类法》的类号，同时应注明《国际十进分类法UDC》的类号；
\item 密级  在右上角注明密级；
\item “博士学位论文”用大号字标明；
\item 题名和副题名   用大号字标明；
\item 作者姓名；
\item 学科专业名称；
\item 研究方向；
\item 导师姓名，职称；
\item 日期包括论文提交日期和答辩日期；
\item 学位授予单位。
\end{enumerate}

\subsection{题名}

题名是以最恰当、最简明的词语反映论文中最重要的特定内容的逻辑组合。

题名所用每一词语必须考虑到有助于选定关键词和编写题录、索引等二次文献可以提供
检索的特定实用信息。

题名应避免使用不常见的缩略词、首字母缩写字、字符、代号和公式等。

题名一般不宜超过20字。

论文应有外文题名，外文题名一般不宜超过10个实词。

可以有副题名。

题名在整本论文中不同地方出现时，应完全相同。

\subsection{前言}

前言是作者对本论文基本特征的简介，如论文背景、主旨、目的、意义等并简述本论文
的创新性成果。

\subsection{摘要}

摘要是论文内容不加注释和评论的简单陈述。

论文应有中、英文摘要，中、英文摘要内容应相同。

摘要应具有独立性和自含性，即不阅读论文的全文，便能获得必要的信息，摘要
中有数据、有结论，是一篇完整的短文，可以独立使用，可以引用，可以用于推广。摘
要的内容应包括与论文同等量的主要信息，供读者确定有无必要阅读全文，也供文摘等
二次文献引用。摘要的重点是成果和结论。

中文摘要一般在1500字，英文摘要不宜超过1500实词。

摘要中不用图、表、化学结构式、非公知公用的符号和术语。

\subsection{关键词}

关键词是为了文献标引工作从论文中选取出来用于表示全文主题内容信息款目的单词或
术语。

每篇论文选取3－8个词作为关键词，以显著的字符另起一行，排在摘要的左下方。在英
文摘要的左下方应标注与中文对应的英文关键词。

\subsection{目次页}

目次页由论文的章、节、附录等的序号、名称和页码组成，另页排在摘要的后面。

\subsection{插图和附表清单}

论文中如图表较多，可以分别列出清单并置于目次页之后。

图的清单应有序号、图题和页码。表的清单应有序号、表题和页码。

符号、标志、缩略词、首字母缩写、计量单位、名词、术语等的注释表符号、标志、缩略词、
首字母缩写、计量单位、名词、术语等的注释说明汇集表，应置于图表清单之后。

\section{主体部分}

\subsection{格式}

主体部分由绪论开始，以结论结束。主体部分必须由另页右页开始。每一章必须另页开
始。全部论文章、节、目的格式和版面安排要划一，层次清楚。

\subsection{序号}

\begin{figure}[htbp]
  \centering
  \includegraphics[width= 0.5\textwidth]{scutlogo.eps}\\
  \caption{测试附录中的插图}\label{fig:appendix1}
\end{figure}

论文的章可以写成：第一章。节及节以下均用阿拉伯数字编排序号，如
1.1，1.1.1等。

论文中的图、表、附注、参考文献、公式、算式等一律用阿拉伯数字分别分章依序连续编排
序号。其标注形式应便于互相区别，一般用下例：图1.2；表2.3；附注1）；文献[4]；式
  (6.3)等。

论文一律用阿拉伯数字连续编页码。页码由首页开始，作为第1页，并为右页另页。封页、
封二、封三和封底不编入页码，应为题名页、前言、目次页等前置部分单独编排页码。页码
必须标注在每页的相同位置，便于识别。

\begin{equation}
    C_i = \frac{2E_i}{k_i(k_i-1)}
\end{equation}

附录依序用大写正体A、B、C、$\cdots$编序号，如：附录A。附录中的图、表、式、参考文
献等另行编序号，与正文分开，也一律用阿拉伯数字编码，但在数码前题以附条序码，如图
A.1；表B.2；式(B.3)；文献[A.5]等。

\subsection{绪论}

绪论（综述）：简要说明研究工作的目的、范围、相关领域的前人工作和知识空白、理
论基础和分析，研究设想、研究方法和实验设计、预期结果和意义等。一般在教科书中
有的知识，在绪论中不必赘述。

绪论的内容应包括论文研究方向相关领域的最新进展、对有关进展和问题的评价、本论
文研究的命题和技术路线等；绪论应表明博士生对研究方向相关的学科领域有系统深入
的了解，论文具有先进性和前沿性；

\begin{problem}
测试定理环境。测试定理环境。测试定理环境。测试定理环境。测试定理环境。测试定理环境。
测试定理环境。测试定理环境。测试定理环境。
\end{problem}

为了反映出作者确已掌握了坚实的基础理论和系统的专门知识，具有开阔的科学视野，对研
究方案作了充分论证，绪论应单独成章，列为第一章，绪论的篇幅应达$1\sim 2$万字，不
得少于$1$万字；绪论引用的文献应在$100$篇以上，其中外文文献不少于$60\%$；引用文献
应按正文中引用的先后排列。

\subsection{正文}

论文的正文是核心部分，占主要篇幅。正文必须实事求是，客观真切，准确完备，合乎
逻辑，层次分明，简便可读。

\begin{figure}[htbp]
  \centering
  \includegraphics[width= 0.5\textwidth]{scutlogo.eps}\\
  \caption{测试附录中的插图}\label{fig:appendix2}
\end{figure}

正文的每一章(除绪论外)应有小结，在小结中应明确阐明作者在本章中所做的工作，特
别是创新性成果。凡本论文要用的基础性内容或他人的成果不应单独成章，也不应作过
多的阐述，一般只引结论、使用条件等，不作推导。

\subsubsection{图}

图包括曲线图、构造图、示意图、图解、框图、流程图、记录图、布置图、地图、照片
、图版等。

图应具有“自明性”，即只看图、图题和图例，不阅读正文，就可以理解图意。

图应编排序号。每一图应有简短确切的图题，连同图号置于图下。必要时，应将图上的
符号、标记、代码，以及实验条件等，用最简练的文字，横排于图题下方，作为图例说
明。

\begin{example}
测试定理环境。测试定理环境。测试定理环境。测试定理环境。测试定理环境。测试定理环境。
测试定理环境。测试定理环境。测试定理环境。
\end{example}

曲线图的纵、横坐标必须标注“量、标准规定符号、单位”。此三者只有在不必要标明
(如无量纲等)的情况下方可省略。坐标上标注的量的符号和缩略词必须与正文一致。

照片图要求主题和主要显示部分的轮廓鲜明，便于制版。如用放大缩小的复制品，必须
清晰，反差适中。照片上应该有表示目的物尺寸的标度。

\subsubsection{表}

表的编排，一般是内容和测试项目由左至右横读，数据依序竖排。表应有自明性。

表应编排序号。

每一表应有简短确切的表题，连同标号置于表上。必要时，应将表中的符号、标记、代
码，以及需要说明事项，以最简练的文字，横排于表题下，作为表注，也可以附注于表
下。表内附注的序号宜用小号阿拉伯数字并加圆括号置于被标注对象的右上角，如：
xxx${}^{1)}$；不宜用“*”，以免与数学上共轭和物质转移的符号相混。

表的各栏均应标明“量或测试项目、标准规定符号、单位”。只有在无必要标注的情况下
方可省略。表中的缩略词和符号，必须与正文中一致。

表内同一栏的数字必须上下对齐。表内不宜用“同上”，“同左”和类似词，一律填入具体数字
或文字。表内“空白”代表未测或无此项，“－”或“\textellipsis”（因“－”可能与代表阴性
  反应相混）代表未发现，“0”代表实测结果确为零。

如数据已绘成曲线图，可不再列表。

\subsubsection{数学、物理和化学式}

正文中的公式、算式或方程式等应编排序号，序号标注于该式所在行(当有续行时，应
标注于最后一行)的最右边。

较长的式，另行居中横排。如式必须转行时，只能在$+$，$-$，$\times$，$\div$，$<$，
$>$处转行。上下式尽可能在等号“$=$”处对齐。

小数点用“$.$”表示。大于$999$的整数和多于三位数的小数，一律用半个阿拉伯数字符的小
间隔分开，不用千位撇。对于纯小数应将$0$列于小数点之前。

示例：应该写成$94\ 652.023\ 567$和$0.314\ 325$, 不应写成$94,652.023,567$和
$.314,325$。

应注意区别各种字符，如：拉丁文、希腊文、俄文、德文花体、草体；罗马数字和阿拉伯数
字；字符的正斜体、黑白体、大小写、上下脚标（特别是多层次，如“三踏步”）、上下偏差
等。

\subsubsection{计量单位}

报告、论文必须采用国务院发布的《中华人民共和国法定计量单位》，并遵照《中华人
民共和国法定计量单位使用方法》执行。使用各种量、单位和符号，必须遵循附录B所
列国家标准的规定执行。单位名称和符号的书写方式一律采用国际通用符号。

\subsubsection{符号和缩略词}

符号和缩略词应遵照国家标准的有关规定执行。如无标准可循，可采纳本学科或本专业
的权威性机构或学术团体所公布的规定；也可以采用全国自然科学名词审定委员会编印
的各学科词汇的用词。如不得不引用某些不是公知公用的、且又不易为同行读者所理解
的、或系作者自定的符号、记号、缩略词、首字母缩写字等时，均应在第一次出现时一
一加以说明，给以明确的定义。

\subsection{结论}

报告、论文的结论是最终的、总体的结论，不是正文中各段的小结的简单重复。结论应
该准确、完整、明确、精炼。在结论中要清楚地阐明论文中有那些自己完成的成果，特
别是创新性成果；

如果不可能导出应有的结论，也可以没有结论而进行必要的讨论。可以在结论或讨论中
提出建议、研究设想、仪器设备改进意见、尚待解决的问题等。

\subsection{致谢}

可以在正文后对下列方面致谢：

\begin{itemize}
\item 国家科学基金、资助研究工作的奖学金基金、合作单位、资助或支持的企业、组织或个
人；
\item 协助完成研究工作和提供便利条件的组织或个人；
\item 在研究工作中提出建议和提供帮助的人；
\item 给予转载和引用权的资料、图片、文献、研究思想和设想的所有者；
\item 其他应感谢的组织或个人。
\end{itemize}

\subsection{参考文献表}

\subsubsection{专著著录格式}

主要责任者，其他责任者，书名，版本，出版地：出版者，出版年

例：1. 刘少奇，论共产党员的修养，修订2版，北京：人民出版社，1962

\subsubsection{连续出版物中析出的文献著录格式}

析出文献责任者，析出文献其他责任者，析出题名，原文献题名，版本：文献中的位置。

例：2. 李四光，地壳构造与地壳运动，中国科学，1973 (4)：400－429

参考文献采用顺序编码制，按论文正文所引用文献出现的先后顺序连续编码。

\section{附录}

附录是作为报告、论文主体的补充项目，并不是必需的。

下列内容可以作为附录编于报告、论文后，也可以另编成册；

\begin{enumerate}
\item 为了整篇论文材料的完整，但编入正文又有损于编排的条理和逻辑性，这一材料
包括比正文更为详尽的信息、研究方法和技术更深入的叙述，建议可以阅读的参考文献
题录，对了解正文内容有用的补充信息等；
\item 由于篇幅过大或取材于复制品而不便于编入正文的材料；
\item 不便于编入正文的罕见珍贵资料；
\item 对一般读者并非必要阅读，但对本专业同行有参考价值的资料；
\item 某些重要的原始数据、数学推导、计算程序、框图、结构图、注释、统计表、计
算机打印输出件等。
\end{enumerate}

附录与正文连续编页码。

每一附录均另页起。

\section{结尾部分 (必要时)}

为了将论文迅速存储入电子计算机，可以提供有关的输入数据。可以编排分类索引、著者索
引、关键词索引等。

% 参考文献。应放在\backmatter之前。
% 推荐使用BibTeX，若不使用BibTeX时注释掉下面一句。
\nocite{*}
\bibliography{main}
% 不使用 BibTeX
%\begin{thebibliography}{2}
%
%\bibitem{deng:01a}
%{邓建松,彭冉冉,陈长松}.
%\newblock {\em \LaTeXe{}科技排版指南}.
%\newblock 科学出版社,书号:7-03-009239-2/TP.1516, 北京, 2001.
%
%\bibitem{wang:00a}
%王磊.
%\newblock {\em \LaTeXe{}插图指南}.
%\newblock 2000.
%\end{thebibliography}

% 附录，必须放在参考文献后，backmatter前
\appendix
\chapter{图论基础知识}
\Blindtext

%%%%%%%%%%%%%%%%%%%%%%%%%%%%%%%%%%%%%%%%%%%%%%%%%%%%%%%%%%%%%%%%%%%%%%%%%%%%%%%
% 书籍附件
\backmatter
%%%%%%%%%%%%%%%%%%%%%%%%%%%%%%%%%%%%%%%%%%%%%%%%%%%%%%%%%%%%%%%%%%%%%%%%%%%%%%%
% 作者简历与科研成果页，应放在backmatter之后
\begin{resume}
% 论文作者身份简介，一句话即可。
\begin{authorinfo}
\noindent 韦小宝，男，汉族，1985年11月出生，江苏省扬州人。
\end{authorinfo}
% 论文作者教育经历列表，按日期从近到远排列，不包括将要申请的学位。
\begin{education}
\item[2007年9月 --- 2010年6月] 南京大学计算机科学与技术系 \hfill 硕士
\item[2003年9月 --- 2007年6月] 南京大学计算机科学与技术系 \hfill 本科
\end{education}
% 论文作者在攻读学位期间所发表的文章的列表，按发表日期从近到远排列。
\begin{publications}
\item Xiaobao Wei, Jinnan Chen, ``Voting-on-Grid Clustering for Secure
  Localization in Wireless Sensor Networks,'' in \textsl{Proc. IEEE International
    Conference on Communications (ICC) 2010}, May. 2010.
\item Xiaobao Wei, Shiba Mao, Jinnan Chen, ``Protecting Source Location Privacy
  in Wireless Sensor Networks with Data Aggregation,'' in \textsl{Proc. 6th
    International Conference on Ubiquitous Intelligence and Computing (UIC)
    2009}, Oct. 2009.
\end{publications}
% 论文作者在攻读学位期间参与的科研课题的列表，按照日期从近到远排列。
\begin{projects}
\item 国家自然科学基金面上项目``无线传感器网络在知识获取过程中的若干安全问题研究''
（课题年限~2010年1月 --- 2012年12月），负责位置相关安全问题的研究。
\item 江苏省知识创新工程重要方向项目下属课题``下一代移动通信安全机制研究''
（课题年限~2010年1月 --- 2010年12月），负责LTE/SAE认证相关的安全问题研究。
\end{projects}
\end{resume}

%%%%%%%%%%%%%%%%%%%%%%%%%%%%%%%%%%%%%%%%%%%%%%%%%%%%%%%%%%%%%%%%%%%%%%%%%%%%%%%
% 生成《学位论文出版授权书》页面，应放在最后一页
% \makelicense

%%%%%%%%%%%%%%%%%%%%%%%%%%%%%%%%%%%%%%%%%%%%%%%%%%%%%%%%%%%%%%%%%%%%%%%%%%%%%%%
\end{document}
